% Compile this UTF-8 document with XeLaTeX or LuaLaTeX.
\documentclass[11pt]{article}
\usepackage{fontspec}
\usepackage{amsmath,amssymb,mathtools,bm}
\usepackage[margin=1in]{geometry}
\setlength{\parindent}{0pt}
\setlength{\parskip}{0.5em}
\begin{document}
\section*{3 Multiple Linear Regression - I}
{\small Source: Regression Modeling with Actuarial and Financial Applications (Frees)/Regression Modeling with Actuarial and Financial Applications.pdf\par}
\medskip
3
 Multiple Linear Regression \(- \mathbf{I}\)
 Chapter Preview. This chapter introduces linear regression in the case of several explana-
 tory variables, known as multiple linear regression. Many basic linear regression concepts
extend directly, including goodness-of-fit measures such as \(R^{2}\) and inference using \(t\)
 statistics. Multiple linear regression models provide a framework for summarizing highly
complex, multivariate data. Because this framework requires only linearity in the param-
eters, we are able to fit models that are nonlinear functions of the explanatory variables,
 thus providing a wide scope of potential applications.
 3.1 Method of Least Squares
 Chapter 2 dealt with the problem of a response depending on a single explanatory
variable. We now extend the focus of that chapter and study how a response may
 depend on several explanatory variables.
@ EMPIRICAL
Filename is Example: Term Life Insurance. Like all firms, life insurance companies con-
"TermLife" tinually seek new ways to deliver products to the market. Those involved in
 product development want to know who buys insurance and how much they buy.
In economics, this is known as the demand side of a market for products. Ana-
lysts can readily get information on characteristics of current customers through
company databases. Potential customers, those who do not have insurance with
 the company, are often the main focus for expanding market share.
In this example, we examine the Survey of Consumer Finances (SCF), a
nationally representative sample that contains extensive information on assets,
liabilities, income, and demographic characteristics of those sampled (potential
U.S. customers). We study a random sample of 500 households with positive
incomes that were interviewed in the 2004 survey. We initially consider the
subset of \(n \,=2 7 5\) families that purchased term life insurance. We address the
second portion of the demand question and determine family characteristics that
influence the amount of insurance purchased. Chapter 11 will consider the first
portion, whether a household purchases insurance, through models where the
 response is a binary random variable.
For term life insurance, the quantity of insurance is measured by the policy
FACE, the amount that the company will pay in the event of the death of the
70
\clearpage
3.1 Method of Least Squares 71
 Table 3.1 Term Life
Standard Summary Statistics
Variable Mean Median Deviation Minimum Maximum
FACE 747,581 150,000 1,674,362 800 14,000,000
INCOME 208,975 65,000 824,010 260 10,000,000
EDUCATION 14.524 16.000 2.549 2.000 17.000
NUMHH 2.960 3.000 1.493 1.000 9.000
LNFACE 11.990 11.918 1.871 6.685 16.455
LNINCOME 11.149 11.082 1.295 5.561 16.118
named insured. Characteristics that will turn out to be important include annual
INCOME; the number of years of EDUCATION of the survey respondent, and
 the number of household members, NUMHH.
In general, we will consider datasets where there are \(k\) explanatory variables
 and one response variable in a sample of size \(n\) . That is, the data consist of
\[
\left\{\begin{aligned} {x_{1 1}} & {{} , \, x_{1 2} , \, \ldots, \, x_{1 k} , \, y_{1}} \\ {x_{2 1}} & {{} , \, x_{2 2} , \, \ldots, \, x_{2 k} , \, y_{2}} \\ {} & {{}} \\ {x_{n 1}} & {{} , \, x_{n 2} , \, \ldots, \, x_{n k} , \, y_{n}} \\ \end{aligned} \right\} .
\]
The \(i\) th observation corresponds to the \(i\) th row, consisting of \(\left( x_{i 1} , \, x_{i \, 2} , \, \ldots, \, x_{i \, k} , \, y_{i} \right)\) .
For this general case, we take \(k+1\) measurements on each entity. For the insur-
ance demand example, \(k=3\) and the data consists of \(( x_{1 1} , \, x_{1 2} , \, x_{1 3} , \, y_{1} ) , \, \ldots\)
\(\left( x_{2 7 5 , 1} , \, x_{2 7 5 , 2} , \, x_{2 7 5 , 3} , \, y_{2 7 5} \right)\) . That is, we use four measurements from each of the
\(n \,=2 7 5\) households.
 Summarizing the Data
We begin the data analysis by examining each variable in isolation of the others. Begin the data
Table 3.1 provides basic summary statistics of the four variables. For FACE and analysis by examining
INCOME, we see that the mean is much greater than the median, suggesting thateach variable in.
the distribution is skewed to the right. Histograms (not reported here) show thatisolation of the others.
this is the case. It will be useful to also consider their logarithmic transforms,
LNFACE and LNINCOME, respectively, which are also reported in Table 3.1.
 The next step is to measure the relationship between each \(x\) on \(y\) , beginning
with the scatter plots in Figure 3.1. The left-hand panel is a plot of FACE versus
INCOME; with this panel, we see a large clustering in the lower-left-hand corner
corresponding to households that have both small incomes and face amounts
of insurance. Both variables have skewed distributions and their joint effect is
 highly nonlinear. The right-hand panel presents the same variables but using
logarithmic transforms. In Figure 3.1, we see a relationship that can be more
readily approximated with a line.
 The term life data are multivariate in the sense that several measurements are
 taken on each household. It is difficult to produce a graph of observations in three
\clearpage
72 Multiple Linear Regression -- I
Table 3.2 Term Life
Correlations NUMHH EDUCATION LNINCOME
EDUCATION \(- 0 . 0 6 4\)
LNINCOME 0.179 0.343
LNFACE 0.288 0.383 0.482
Figure 3.1 Income FACE (n m \(\circ\) \(\check{\circ}\) B 809 \(\omicron^{\mathrm{O}}\)
versus face amount of 14 16 \(\circ\)
 term life insurance. 12
The left panel is a plot 10 14 : \(\vartheta_{\infty}\) 0
 of face versus income, 8 0 12 \(\omicron^{\circ}\) \(\mathcal{S}\) 0 Oa
 showing a highly
 nonlinear pattern. In 6 \(\diamond^{\textsc{c}}\) 10 0 00
 the right-hand panel, 4 H \(\oint\) \(\circ\)
face versus income is
 in natural logarithmic 2
units, suggesting a 0 0
linear (although
variable) pattern. 01 2 3 5 6 8 6 810121416
INCOME (in millions) LNINCOME
 or more dimensions on a two-dimensional platform, such as a piece of paper,
 that is not confusing, misleading, or both. To summarize graphically multivariate
data in regression applications, consider using a scatterplot matrix such as in
Figure 3.2. Each square of this figure represents a simple plot of one variable
versus another. For each square, the row variable gives the units of the vertical
 axis and the column variable gives the units of the horizontal axis. The matrix
is sometimes called a half scatterplot matrix because only the lower-left-hand
elements are presented.
The scatterplot matrix can be numerically summarized using a correlation
matrix. Each correlation in Table 3.2 corresponds to a square of the scatterplot
 matrix in Figure 3.2. Analysts often present tables of correlations because they
are easy to interpret. However, remember that a correlation coefficient merely
measures the extent of linear relationships. Thus, a table of correlations provides
 a sense of linear relationships but may miss a nonlinear relationship that can be
revealed in a scatterplot matrix.
 The scatterplot matrix and corresponding correlation matrix are useful devices
for summarizing multivariate data. They are easy to produce and to interpret. Still,
each device captures only relationships between pairs of variables and cannot
 quantify relationships among several variables.
 Method of Least Squares
Consider this question: can knowledge of education, household size, and income
 help us understand the demand for insurance? The correlations in Table 3.2 and
the graphs in Figures 3.1 and 3.2 suggest that each variable, EDUCATION,
NUMHH, and LNINCOME, may be a useful explanatory variable of LNFACE
\clearpage
3.1 Method of Least Squares 73
Figure 3.2
Scatterplot matrix of
four variables. Each
6 square is a scatter plot.
NUMHH
10 1E EDUCATION
:
:
0
16
12 1 LNINCOME
10 :
6 :
16 LIE 888 D LNFACE 12
1 14
14
12
10
:
:
0 0
2 4 6 8 5 10 15 6810121416 81012 14 16
when taken individually. It seems reasonable to investigate the joint effect of
 these variables on a response.
The geometric concept of a plane is used to explore the linear relationship
 between a response and several explanatory variables. Recall that a plane extends
the concept of a line to more than two dimensions. \(\mathrm{A}\) plane may be defined
 through an algebraic equation such as
\[
y \,=\, b_{0} \,+\, b_{1} x_{1} \,+\, \cdots\,+\, b_{k} \, x_{k} .
\]
 This equation defines a plane in \(k+1\) dimensions. Figure 3.3 shows a plane
in three dimensions. For this figure, there is one response variable, LNFACE,
and two explanatory variables, EDUCATION and LNINCOME (NUMHH is
held fixed). It is difficult to graph more than three dimensions in a meaningful
way.
We need a way to determine a plane from the data. The difficulty is that in
 most regression analysis applications, the number of observations, \(n\) , far exceeds
 the number of observations required to fit a plane, \(k+1\) . Thus, it is generally
 not possible to find a single plane that passes through all \(n\) observations. As in
Chapter \(2 ,\) we use the method of least squares to determine a plane from the data.
\clearpage
74 Multiple Linear Regression -- I
E
\[
\sum_{n=1}^{1 2} \left[ \begin{matrix} {\sqrt{2}} \\ {\sqrt{2}} \\ {\sqrt{2}} \\ {\sqrt{2}} \\ \end{matrix} \right] \left\{\begin{matrix} \sqrt{2} \\ {\sqrt{2}} \\ {\sqrt{2}} \\ {\sqrt{2}} \\ {\sqrt{2}} \\ {\sqrt{2}} \\ {\sqrt{2}} \\ \end{matrix} \right] \left( \begin{matrix} \end{matrix} \right) \left( \begin{matrix} \sqrt{2} \\ {\sqrt{2}} \\ {\sqrt{2}} \\ {\sqrt{2}} \\ {\sqrt{2}} \\ {\sqrt{2}} \\ {\sqrt{2}} \\ {\sqrt{2}} \\ \end{matrix} \right) \left( \begin{matrix} \end{matrix} \right) \left( \begin{matrix} \end{matrix} \right) 1} \\ {\left\{\begin{matrix} \left( \begin{matrix} \end{matrix} \left( \begin{matrix} \end{matrix} \right) \left( \left( \begin{matrix} \end{matrix} \left( \left( \begin{matrix} \end{matrix} \end{matrix} \right) \right) \right) \right)} \\ {\left\{\begin{matrix} \left( \left( \left( \begin{matrix} \end{matrix} \right) \left( \left( \left( \left( \left( \left( \left( \left( \left( \left1 \right) \right) \right) \right) \right) \right) \right) \right) \right)} \\ \end{matrix} \right)} \\ {\left\{\begin{matrix} \left( \begin{matrix} \end{matrix} \right) \left( \begin{matrix} \left( \left( \begin{matrix} \end{matrix} \right) \right) \right) \left( \left( \begin{matrix} \left( \left( \begin{matrix} \end{matrix} \right) \right) 1} \\ \end{matrix} \right) \right. ( \begin{matrix} \end{matrix} \right) 1 \right. \right. \begin{matrix} \end{matrix} \right) ( \begin{matrix} \end{matrix} \right) \right.
\]
 example of a
three-dimensional
The method of least squares is based on determining the values of \(b_{0}^{*}\)
\(b_{1}^{*} , \dots, b_{k}^{*}\) that minimize the quantity
\[
S S \big( b_{0}^{*} , b_{1}^{*} , \dots, b_{k}^{*} \big)=\sum_{i=1}^{n} \left( y_{i}-\left( b_{0}^{*}+b_{1}^{*} x_{i 1}+\dots+b_{k}^{*} x_{i k} \right) \right)^{2} . \tag{3.1}
\]
We drop the asterisk notation and use \({b_{0}} , {b_{1}} , \dots, {b_{k}}\) to denote the best values,
 known as the least squares estimates. With the least squares estimates, define the
 least squares, or fited , regression plane as
\[
{\widehat y} \,=b_{0} \,+\, b_{1} x_{1} \,+\, \cdots\,+\, b_{k} x_{k} .
\]
The least squares estimates are determined by minimizing \(\SS( b_{0}^{\ast} , \, b_{1}^{\ast} , \ldots, \, b_{k}^{\ast} )\) It
is difficult to write down the resulting least squares estimators using a simple
formula unless one resorts to matrix notation. Because of their importance in
applied statistical models, an explicit formula for the estimators is provided
here. However, these formulas have been programmed into a wide variety of
statistical and spreadsheet software packages. The fact that these packages are
 readily available allows data analysts to concentrate on the ideas of the estimation
 procedure instead of focusing on the details of the calculation procedures.
As an example, a regression plane was fit to the term life data where three
 explanatory variables, \(x_{1}\) for EDUCATION, \(x_{2}\) for NUMHH, and \(x_{3}\) for LNIN-
COME, were used. The resulting fitted regression plane is
\[
\widehat{y} \,=\, 2 . 5 8 4 \,+\, 0 . 2 0 6 x_{1} \,+\, 0 . 3 0 6 x_{2} \,+\, 0 . 4 9 4 x_{3} . \tag{3.2}
\]
 Matrix Notation
 Assume that the data are of the form \(\left( x_{i \, 0} , \, x_{i \, 1} , \, \ldots, \, x_{i \, k} , \, y_{i} \right)\) where \(i=1 , \dots, n\) .
Here, the variable \(x_{i \, 0}\) is associated with the " intercept"term. That is, in most
applications, we assume that \(x_{i 0}\) is identically equal to \(1\) and thus need not be
explicitly represented. However, there are important applications where this is
not the case; thus, to express the model in general notation, it is included here.
\clearpage
3.1 Method of Least Squares 75
 The data are represented in matrix notation using
\[
\mathbf{y}=\left( \begin{matrix} {y_{1}} \\ {y_{2}} \\ {\vdots} \\ {y_{n}} \\ \end{matrix} \right) \ \ \ \mathrm{a n d} \ \ \ \mathbf{X}=\left( \begin{matrix} {x_{1 0}} & {x_{1 1}} & {\cdots} & {x_{1 k}} \\ {x_{2 0}} & {x_{2 1}} & {\cdots} & {x_{2 k}} \\ {\vdots} & {\vdots} & {\ddots} & {\vdots} \\ {x_{n 0}} & {x_{n 1}} & {\cdots} & {x_{n k}} \\ \end{matrix} \right) .
\]
Here, \({\bf y}\) is the \(n \times1\) vector of responses and \({\bf X}\) is the \(n \times( k+1 )\) matrix of
explanatory variables. We use the matrix algebra convention that lower- and
uppercase bold letters represent vectors and matrices, respectively. (If you need
 to brush up on matrices, review Section 2.11.)
 Example: Term Life Insurance, Continued. Recall that \(y\) represents the loga-
rithmic face; \(x_{1}\) , years of education; \(x_{2}\) , number of household members; and \(x_{3}\) :
logarithmic income. Thus, there are \(k=3\) explanatory variables and \(n \,=2 7 5\)
 households. The vector of responses and the matrix of explanatory variables are
\[
\mathbf{y}=\left( \begin{matrix} {y_{1}} \\ {y_{2}} \\ {\vdots} \\ {y_{2 7 5}} \\ \end{matrix} \right)=\left( \begin{matrix} {9 . 9 0 4} \\ {1 1 . 7 7 5} \\ {\vdots} \\ {9 . 2 1 0} \\ \end{matrix} \right) \quad\mathrm{a n d}
\]
\[
\mathbf{X}=\left( \begin{matrix} {1} & {x_{1 1}} & {x_{1 2}} & {x_{1 3}} \\ {1} & {x_{2 1}} & {x_{2 2}} & {x_{2 3}} \\ {\vdots} & {\vdots} & {\vdots} & {\vdots} \\ {1} & {x_{2 7 5 , 1}} & {x_{2 7 5 , 2}} & {x_{2 7 5 , 3}} \\ \end{matrix} \right)=\left( \begin{matrix} {1} & {1 6} & {3} & {1 0 . 6 6 9} \\ {1} & {9} & {3} & {9 . 3 9 3} \\ {\vdots} & {\vdots} & {\vdots} & {\vdots} \\ {1} & {1 2} & {1} & {1 0 . 5 4 5} \\ \end{matrix} \right) .
\]
For example, for the first observation in the dataset, the dependent variable is
\(y_{1}=9 . 9 0 4\) (corresponding to \(\operatorname{e x p} ( 9 . 9 0 4 )=\S2 0 , 0 0 0 )\) for a survey respondent
with 16 years of education living in a household with 3 people with logarithmic
income of 10.669 (exp(10.669) = \$43,000).
Under the least squares estimation principle, our goal is to choose the coeff-
cients \(b_{0}^{\ast} , \, b_{1}^{\ast} , \dots, \, b_{k}^{\ast}\) to minimize the sum of squares function \(\SS( b_{0}^{\ast} , \, b_{1}^{\ast} , \ldots, \, b_{k}^{\ast} )\) .
Using calculus, we return to equation (3.1), take partial derivatives with respect
to each coefficient, and set these quantities equal to zero:
\[
\begin{aligned} {\frac{\partial} {\partial b_{j}^{*}} S S \big( b_{0}^{*} , b_{1}^{*} , \dots, b_{k}^{*} \big)} & {{}=\sum_{i=1}^{n} (-2 x_{i j} ) \left( y_{i}-\left( b_{0}^{*}+b_{1}^{*} x_{i 1}+\dots+b_{k}^{*} x_{i k} \right) \right)} \\ {} & {{}=0 , \; \; \; \mathrm{f o r} \; \; j=0 , 1 , \dots, k .} \\ \end{aligned}
\]
This is a system of \(k+1\) equations and \(k+1\) unknowns that can be readily
 solved using matrix notation, as follows.
We may express the vector of parameters to be minimized as \({\bf b}^{*}=( b_{0}^{*}\) .
\(b_{1}^{*} , \dots, b_{k}^{*} )^{\prime}\) Using this, the sum of squares can be written as \(S S ( {\bf b^{*}} )=( {\bf y}-{\bf X b^{*}} )^{\prime}\)
\clearpage
76 Multiple Linear Regression -- I
\(( \mathbf{y}-\mathbf{X b}^{*} )\) . Thus, in matrix form, the solution to the minimization problem can
be expressed as \(( \partial/ \partial{\bf b}^{*} ) S S ( {\bf b}^{*} )={\bf0}\) . This solution satisfies the normal equations
\[
\mathbf{X}^{\prime} \mathbf{X} \mathbf{b}=\mathbf{X}^{\prime} \mathbf{y} . \tag{3.3}
\]
Here, the asterisk notation (*) has been dropped to denote the fact that \(\mathbf{b}=\)
\(( b_{0} , \, b_{1} , \, \dots, \, b_{k} )^{\prime}\) represents the best vector of values in the sense of minimizing
\(S S ( \mathbf{b^{*}} )\) over all choices of \(\mathbf{b^{*}}\) .
 The least squares estimator \(\mathbf{b}\) need not be unique. However, assuming that the
 explanatory variables are not linear combinations of one another, we have that
\(\mathbf{X}^{\prime} \mathbf{X}\) is invertible. In this case, we can write the unique solution as
\[
\mathbf{b}=\left( \mathbf{X}^{\prime} \mathbf{X} \right)^{-1} \mathbf{X}^{\prime} \mathbf{y} . \tag{3.4}
\]
To illustrate, for the term life example, equation (3.2) yields
\[
\mathbf{b}=\left( \begin{matrix} {b_{0}} \\ {b_{1}} \\ {b_{2}} \\ {b_{3}} \\ \end{matrix} \right)=\left( \begin{matrix} {2 . 5 8 4} \\ {0 . 2 0 6} \\ {0 . 3 0 6} \\ {0 . 4 9 4} \\ \end{matrix} \right) .
\]
 3.2 Linear Regression Model and Properties of Estimators
In the previous section, we learned how to use the method of least squares to
fit a regression plane with a dataset. This section describes the assumptions
underpinning the regression model and some of the resulting properties of the
regression coefficient estimators. With the model and the fitted data, we will be
 able to draw inferences about the sample dataset to a larger population. Moreover,
we will later use the regression model assumptions to help us improve the model
 specification in Chapter \(5\) .
3.2.1 Regression Function
 Most of the assumptions of the multiple linear regression model carry over directly
from the basic linear regression model assumptions introduced in Section 2.2.
The primary difference is that we now summarize the relationship between the
response and the explanatory variables through the regression function
\[
\textrm{E} y=\beta_{0} x_{0}+\beta_{1} x_{1}+\cdots+\beta_{k} x_{k} , \tag{3.5}
\]
which is linear in the parameters \(\beta_{0} , \ldots, \beta_{k} .\) Henceforth, we will use \(x_{0}=1\) for
 the variable associated with the parameter \(\beta_{0}\) ; this is the default in most statistical
 packages, and most applications of regression include the intercept term \(\beta_{0}\) The
Interpret \(\beta_{\mathrm{j}}\) to be the intercept is the expected value of \(y\) when all of the explanatory variables are
expected change in \({\bf y}\) equal to zero. Although rarely of interest, the term \(\beta_{0}\) serves to set the height of
 per unit change in \(\mathrm{x_{j}}\) the fitted regression plane.
 assuming all other In contrast, the other betas are typically important parameters from a regression
explanatory variables
are held fred. study. To help interpret them, we initially assume that \(x_{j}\) varies continuously and
\clearpage
 3.2 Linear Regression Model and Properties of Estimators 77
is not related to the other explanatory variables. Then, we can interpret \(\beta_{j}\) as
the expected change in \(y\) per unit change in \(x_{j}\) , assuming all other explanatory
variables are held fixed. That is, from calculus, you will recognize that \(\beta_{j}\) can be
interpreted as a partial derivative. Specifically, using equation (3.5), we have
\[
\beta_{j}=\frac{\partial} {\partial x_{j}} \mathrm{E} ~ y .
\]
3.2.2 Regression Coefficient Interpretation
Let us examine the regression coefficient estimates from the term life insurance
example and focus initially on the \(s i g n\) of the coefficients. For example, from
equation (3.2), the coefficient associated with NUMHH is \(b_{2}=0 . 3 0 6 > 0 .\) If
we consider two households that have the same income and the same level
of education, then the larger household (in terms of NUMHH) is expected to
 demand more term life insurance under the regression model. This is a sensible
 interpretation, as larger households have more dependents for which term life
 insurance can provide needed financial assets in the event of the untimely death
of a breadwinner. The positive coefficient associated with income \(b_{3}=0 . 4 9 4 )\)
 is also plausible; households with larger incomes have more disposable dollars to
 purchase insurance. The positive sign associated with EDUCATION ( \(b_{1}=0 . 2 0 6 )\)
is also reasonable; more education suggests that respondents are more aware of
their insurance needs, other things being equal.
You will also need to interpret the amount of the regression coefficient.
Consider first the EDUCATION coefficient. Using equation (3.2), fitted val-
ues of \(\widehat{L N F A C E}\) were calculated by allowing EDUCATION to vary and keeping
NUMHH and LNINCOME fixed at the sample averages. The results are as
follows:
Effects of Small Changes in Education
EDUCATION 14 14.1 14.2 14.3
\(\widehat{L N F A C E}\) 11.883 11.904 11.924 11.945
\(\widehat{F A C E}\) 144,803 147,817 150,893 154,034
FACE \% Change 2.081 2.081 2.081
 As EDUCATION increases, \(\widehat{L N F A C E}\) increases. Further, the amount of \(\widehat{L N F A C E}\)
increase is a constant 0.0206. This comes directly from equation (3.2); as EDU-
CATION increases by 0.1 years, we expect the demand for insurance to increase
by 0.0206 logarithmic dollars, holding NUMHH and LNINCOME fixed. This
interpretation is correct, but most product development directors are not overly
 fond of logarithmic dollars. To return to dollars, fitted face values can be calculated
through exponentiation as \(\widehat{F A C E}=\operatorname{e x p} ( \widehat{L N F A C E} )\) Moreover, the percentage
change can be computed; for example, \(1 0 0 \times( 1 4 7 , 8 1 7 / 1 4 4 , 8 0 3-1 ) \approx2 . 0 8 7 o .\)
\clearpage
78 Multiple Linear Regression -- I
This provides another interpretation of the regression coefficient; as EDUCA-
TION increases by 0.1 years, we expect the demand for insurance to increase by
\(2 . 0 8 7_{o}\) . This is a simple consequence of calculus, using 3ln \(y / \partial x \,=\left( \partial y / \partial x \right) / y\) :
that is, a small change in the logarithmic value of \(y\) equals a small change in \(y\)
 as a proportion of \(y\) . It is because of this calculus result that we use natural logs
 instead of common logs in regression analysis. Because this table uses a discrete
change in EDUCATION, the 2.08\% differs slightly from the continuous result
\(0 . 2 0 6 \times\) (change in EDUCATION) = 2.06\%. However, this proximity is usually
regarded as suitable for interpretation purposes.
Continuing this logic, consider small changes in logarithmic income.
Effects of Small Changes in Logarithmic Income
LNINCOME 11 11.1 11.2 11.3
INCOME 59,874 66,171 73,130 80,822
INCOME \(\gamma_{o}\) Change 10.52 10.52 10.52
LNFACE 11.957 12.006 12.055 12.105
\(\widehat{F A C E}\) 155,831 163,722 172,013 180,724
\(\widehat{F A C E}\) \% Change 5.06 5.06 5.06
\(\widehat{F A C E}\) \% Change/INCOME \(\gamma_{o}\) Change 0.482 0.482 0.482
We can use the same logic to interpret the LNINCOME coefficient in equation
(3.2). As logarithmic income increases by 0.1 units, we expect the demand for
insurance to increase by 5.06\%. This takes care of logarithmic units in the \(y\)
but not the \(x\) . We can use the same logic to say that as logarithmic income
increases by 0.1 units, INCOME increases by 10.52\%. Thus, a 10.52\% change
in INCOME corresponds to a 5.06\% change in FACE. To summarize, we say
that, holding NUMHH and EDUCATION fixed, we expect that a \(1 7 o\) increase in
INCOME is associated with a 0.482\% increase in \(\widehat{F A C E}\) (as earlier, this is close
 to the parameter estimate \(b_{3}=0 . 4 9 4 )\) . The coefficient associated with income is
known as an elasticity in economics. In economics, elasticity is the ratio of the
 percentage change in one variable to the percentage change in another variable.
Mathematically, we summarize this as
\[
\frac{\partial\, \operatorname{l n} \, y} {\partial\, \operatorname{l n} \, x}=\left( \frac{\partial\, \, y} {y} \right) \bigg/ \bigg( \frac{\partial\, \, x} {x} \bigg) .
\]
 3.2.3 Model Assumptions
As in Section 2.2 for a single explanatory variable, there are two sets of assump-
 tions that one can use for multiple linear regression. They are equivalents sets,
each having comparative advantages as we proceed in our study of regression.
 The observables representation focuses on variables of interest \(( x_{i \, 1} , \, \dots, \, x_{i \, k} , \, y_{i} )\) .
\clearpage
 3.2 Linear Regression Model and Properties of Estimators 79
 The error representation provides a springboard for motivating our goodness-of-
fit measures and study of residual analysis. However, the latter set of assumptions
 focuses on the additive errors case and obscures the sampling basis of the model.
 Multiple Linear Regression Model Sampling Assumptions
 Observables Representation Error Representation
F1.E \(y_{i} \,=\, \beta_{0} \,+\, \beta_{1} x_{i 1} \,+\, \cdots\,+\, \beta_{k} x_{i k} .\) E1. \(y_{i} \,=\, \beta_{0} \,+\, \beta_{1} x_{i 1} \,+\, \cdots\,+\, \beta_{k} x_{i k} \,+\, \varepsilon_{i} \, .\)
F2. \(\{x_{i 1} , \, \dots, \, x_{i k} \}\) E2. \(\{x_{i 1} , \, \dots, \, x_{i k} \}\)
 are nonstochastic variables. are nonstochastic variables.
F3. Var \(y_{i}=\sigma^{2}\) . E3.E \(\varepsilon_{i}=0\) and Var \(\varepsilon_{i}=\sigma^{2}\) .
F4. \(\{y_{i} \}\) are independent random variables. E4. \(\{\varepsilon_{i} \}\) are independent random variables.
F5. \(\{y_{i} \}\) are normally distributed. E5. \(\{\varepsilon_{i} \}\) are normally distributed.
 To further motivate Assumptions F2 and \(\mathrm{F 4}\) we usually assume that our data
have been realized as the result of a stratified sampling scheme, where each
 unique value of \(\{x_{i \, 1} , \, \dots, \, x_{i \, k} \}\) is treated as a stratum. That is, for each value of
\(\{x_{i 1} , \, \dots, \, x_{i k} \}\) , we draw a random sample of responses from a population. Thus,
responses in each stratum are independent from one another, as are responses
 from different strata. Chapter \(6\) will discuss this sampling basis in further detail.
3.2.4 Properties of Regression Coefficient Estimators
Section 3.1 described the least squares method for estimating regression coef-
ficients. With the regression model assumptions, we can establish some basic
properties of these estimators. To do this, from Section 2.11.4, we have that the
expectation of a vector is the vector of expectations, so that
\[
\textrm{E}_{\mathbf{y}}=\left( \begin{matrix} {\textrm{E}_{1}} \\ {\textrm{E}_{2}} \\ {\vdots_{\phantom{\Big|}}} \\ {\textrm{E}_{n}} \\ \end{matrix} \right) .
\]
Further, basic matrix multiplication shows that
\[
\mathbf{X} \boldsymbol{\beta}=\left( \begin{matrix} {1} & {x_{1 1}} & {\cdots} & {x_{1 k}} \\ {1} & {x_{2 1}} & {\cdots} & {x_{2 k}} \\ {\vdots} & {\vdots} & {\ddots} & {\vdots} \\ {1} & {x_{n , 1}} & {\cdots} & {x_{n , k}} \\ \end{matrix} \right) \left( \begin{matrix} {\beta_{0}} \\ {\beta_{1}} \\ {\vdots} \\ {\beta_{k}} \\ \end{matrix} \right)=\left( \begin{matrix} {\beta_{0}+\beta_{1} x_{1 1}+\cdots+\beta_{k} x_{1 k}} \\ {\beta_{0}+\beta_{1} x_{2 1}+\cdots+\beta_{k} x_{2 k}} \\ {\vdots} \\ {\beta_{0}+\beta_{1} x_{n 1}+\cdots+\beta_{k} x_{n k}} \\ \end{matrix} \right) .
\]
Because the \(i\) th row of assumption F1 is E \(y_{i}=\beta_{0}+\beta_{1} x_{i 1}+\cdots+\beta_{k} x_{i k}\) , we
can rewrite this assumption in matrix formulation as \(\textrm{E}_{\mathbf{y}}=\mathbf{X} \boldsymbol{\beta}\) . We are now in a
 position to state the first important property of least squares regression estimators.
Property 1. Consider a regression model and let Assumptions F1--F4hold.
Then, the estimator \(\mathbf{b}\) defined in equation (3.4) is an unbiased estimator of the
 parameter vector \(\boldsymbol{\beta}\) .
\clearpage
80 Multiple Linear Regression -- I
To establish Property 1, we have that
\[
\mathrm{\boldmath~ E ~ b}=\mathrm{\boldmath~ E ~} \left( {\left( {{\bf X^{\prime} X}} \right)^{-1}} {\bf X^{\prime}} \right)={\left( {{\bf X^{\prime} X}} \right)^{-1}} {\bf X^{\prime}} \mathrm{\boldmath~ E ~} \mathbf y={\left( {{\bf X^{\prime} X}} \right)^{-1}} {\bf X^{\prime}} \left( {{\bf X \beta}} \right)=\beta,
\]
using matrix multiplication rules. This chapter assumes that \(\mathbf{X}^{\prime} \mathbf{X}\) is invertible.
One can also show that the least squares estimator need only be a solution of
 the normal equations for unbiasedness (not requiring that \(\mathbf{X}^{\prime} \mathbf{X}\) be invertible, see
Section 4.7.3). Thus, \(\mathbf{b}\) is said to be an unbiased estimator of \(\boldsymbol{\beta}\) . In particular,
E \(b_{j}=\beta_{j}\) for \(j=0 , 1 , \dots, k\) .
 Because independence implies zero covariance, from Assumption F4 we have
\(\mathrm{C o v} ( y_{i} , \, y_{j} )=0\) for \(i \neq j\) From this, Assumption F3 and the definition of the
variance of a vector, we have
\[
\begin{aligned} {\mathrm{V a r ~} \mathbf{y}} & {{}=\left( \begin{matrix} {\mathrm{V a r ~} y_{1}} & {\mathrm{C o v} ( y_{1} , y_{2} )} & {\cdots} & {\mathrm{C o v} ( y_{1} , y_{n} )} \\ {\mathrm{C o v} ( y_{2} , \, y_{1} )} & {\mathrm{V a r ~} y_{2}} & {\cdots} & {\mathrm{C o v} ( y_{2} , \, y_{n} )} \\ {\vdots} & {\vdots} & {\ddots} & {\vdots} \\ {\mathrm{C o v} ( y_{n} , \, y_{1} )} & {\mathrm{C o v} ( y_{n} , \, y_{2} )} & {\cdots} & {\mathrm{V a r ~} y_{n}} \\ \end{matrix} \right)} \\ {} & {{}=\left( \begin{matrix} {\sigma^{2}} & {0} & {\cdots} & {0} \\ {0} & {\sigma^{2}} & {\cdots} & {0} \\ {\vdots} & {\vdots} & {\ddots} & {\vdots} \\ {0} & {0} & {\cdots} & {\sigma^{2}} \\ \end{matrix} \right)=\sigma^{2} \mathbf{I} .} \\ \end{aligned}
\]
where \(\textbf{I}\) is an an \(n \times n\) identity matrix. We are now in a position to state the
second important property of least squares regression estimators.
Property \({\bf2}\) . Consider a regression model and let Assumptions F1-F4 hold.
Then, the estimator \(\textbf{b}\) defined in equation (3.4) has variance Var \(\mathbf{b}=\)
\(\sigma^{2} {( \mathbf{X}^{\prime} \mathbf{X} )}^{-1}\)
To establish Property 2, we have
\[
\begin{aligned} {\mathrm{V a r ~ b}} & {{}=\mathrm{V a r} \left( {{( {\bf X}^{\prime} {\bf X} )}^{-1}} {\bf X}^{\prime} {\bf y} \right)=\left[ {( {\bf X}^{\prime} {\bf X} )}^{-1} {\bf X}^{\prime} \right] \mathrm{V a r} \left( {\bf y} \right) \left[ {\bf X} {{( {\bf X}^{\prime} {\bf X} )}^{-1}} \right]} \\ {} & {{}=\left[ {{( {\bf X}^{\prime} {\bf X} )}^{-1}} {\bf X}^{\prime} \right] \sigma^{2} {\bf I} \left[ {\bf X} {{( {\bf X}^{\prime} {\bf X} )}^{-1}} \right]=\sigma^{2} {{( {\bf X}^{\prime} {\bf X} )}^{-1}} {\bf X}^{\prime} {{\bf X} ( {{\bf X}^{\prime} {\bf X} )}^{-1}}=\sigma^{2} {( {\bf X}^{\prime} {\bf X} )}^{-1}} \\ \end{aligned}
\]
1
 as required. This important property will allow us to measure the precision of the
estimator \(\mathbf{b}\) when we discuss statistical inference. Specifically, by the definition
 of the variance of a vector (see Section 2.11.4),
\[
\mathrm{V a r ~ b}=\left( \begin{matrix} {\mathrm{V a r} \; b_{0}} & {\mathrm{C o v} ( b_{0} , \, b_{1} )} & {\cdots} & {\mathrm{C o v} ( b_{0} , \, b_{k} )} \\ {\mathrm{C o v} ( b_{1} , \, b_{0} )} & {\mathrm{V a r} \; b_{1}} & {\cdots} & {\mathrm{C o v} ( b_{1} , \, b_{k} )} \\ {\vdots} & {\vdots} & {\ddots} & {\vdots} \\ {\mathrm{C o v} ( b_{k} , \, b_{0} )} & {\mathrm{C o v} ( b_{k} , \, b_{1} )} & {\cdots} & {\mathrm{V a r} \; b_{k}} \\ \end{matrix} \right)={\boldsymbol\sigma}^{2} {( {\bf X}^{\prime} {\bf X} )}^{-1} .
\]
.6)
Thus, for example, Var \(b_{j}\) is \(\sigma^{2}\) times the \(( j+1 )\) st diagonal entry of \(\left( \mathbf{X}^{\prime} \mathbf{X} \right)^{-1}\) .
As another example, \(\mathrm{C o v} ( b_{0} , b_{j} )\) is \(\sigma^{2}\) times the element in the first row and
\(( j+1 )\) st column of (X'X)-1
\clearpage
3.3 Estimation and Goodness of Fit 81
Although alternative methods are available that are preferable for specific
 applications, the least squares estimators have proved effective for many routine
 data analyses. One desirable characteristic of least squares regression estimators
 is summarized in the following well-known result.
Gauss-Markov Theorem. Consider the regression model and let Assump-
 tions F1-- F4hold. Then, in the class of estimators that are linear functions of
the responses, the least squares estimator \(\mathbf{b}\) defined in equation (3.4) is the
 minimum variance unbiased estimator of the parameter vector \(\boldsymbol{\beta}\) .
We have already seen in Property 1 that the least squares estimators are The Gauss-Markov
unbiased. The Gauss-Markov theorem states that the least squares estimator is theorem states that the
 the most precise in that it has the smallest variance. (In a matrix context, minimumleast squares.
estimator is the most
 variance means that if \(\mathbf{b^{*}}\) is any other estimator, then the difference of the variance precise in that it has
matrices, Var \(\mathbf{b^{*} \mathrm{-} V a r} ~ \mathbf{b}\) , is nonnegative definite.) the smallest variance.
An additional important property concerns the distribution of the least squares
regression estimators.
Property 3. Consider a regression model and let Assumptions F1-F5 hold.
Then, the least squares estimator \(\mathbf{b}\) defined in equation (3.4) is normally
distributed.
To establish Property 3, we define the weight vectors, \(\mathbf{w}_{i} \,=\, {( \mathbf{X}^{\prime} \mathbf{X} )}^{-1} ( 1\)
\(x_{i 1} , \dots, \, x_{i k} )^{\prime}\) With this notation, we note that 
\[
\mathbf{b=} \mathbf{( X^{\prime} X )}^{-1} \mathbf{X}^{\prime} \mathbf{y}=\sum_{i=1}^{n} \mathbf{w}_{i} \, y_{i} \, ,
\]
 so that \(\mathbf{b}\) is a linear combination of responses. With Assumption F5, the responses
 are normally distributed. Because linear combinations of normally distributed
random variables are normally distributed, we have the conclusion of Property 3.
This result underpins much of the statistical inference that will be presented in
Sections 3.4 and 4.2.
3.3 Estimation and Goodness of Fit
 Residual Standard Deviation
Additional properties of the regression coefficient estimators will be discussed
when we focus on statistical inference. We now continue our estimation discus-
 sion by providing an estimator of the other parameter in the linear regression
 model, \(\sigma^{2}\) .
 Our estimator for \(\sigma^{2}\) can be developed using the principle of replacing theoret-
 ical expectations by sample averages. In examining \(\sigma^{2}=\operatorname{E} \left( y-\operatorname{E} y \right)^{2}\) , replac-
ing the outer expectation by a sample average suggests using the estimator
\(n^{-1} \sum_{i \mathop=1}^{n} ( y_{i} \textrm{-E} y_{i} )^{2}\) . Because we do not observe E \(y_{i}=\beta_{0}+\beta_{1} x_{i 1}+\cdots+\)
\clearpage
82 Multiple Linear Regression -- I
\(\beta_{k} x_{i k}\) , we use in its place the corresponding observed quantity \(b_{0}+b_{1} x_{i 1}\)
\({+\cdots+b_{k} x_{i \, k}=\widehat{y}_{i}}\) This leads to the following:
Defnition . An estimator of \(\sigma^{2}\) , the mean square error (MSE), is defined as
\[
\begin{aligned} {s^{2}=\frac{1} {n-( k+1 )} \sum_{i=1}^{n} \left( y_{i} \,-\, \widehat{y_{i}} \right)^{2} .} \\ \end{aligned} \tag{3.7}
\]
The positive square root, \(s=\sqrt{s^{2}}\) is called the residual standard deviation.
This expression generalizes the definition in equation (2.3), which is valid for
\(k=1\) It turns out that, by using \(n-( k+1 )\) instead of \(n\) in the denominator of
equation (3.7), \(s^{2}\) is an unbiased estimator of \(\sigma^{2}\) Essentially, by using \(\widehat{y_{i}}\) instead
of \(\textrm{E} y_{i}\) in the definition, we have introduced some small dependencies among
the deviations from the responses \(y_{i} \,-\, \widehat{y}_{i}\) , thus reducing the overall variability.
To compensate for this lower variability, we also reduce the denominator in the
 definition of \(s^{2}\) .
To provide further intuition on the choice of \(n-( k+1 )\) in the definition
of \(s^{2}\) , we introduced the concept of residuals in the context of multiple linear
regression. From Assumption E1, recall that the random errors can be expressed
as \(\varepsilon_{i} \,=\, y_{i} \,-( \beta_{0}+\beta_{1} x_{i 1}+\cdots+\beta_{k} x_{i k} ) .\) Because the parameters \(\beta_{0} , \ldots, \beta_{k}\) are
not observed, the errors themselves are not observed. Instead, we examine the
" estimatederrors," or residuals, defined by \(e_{i} \,=\, y_{i} \,-\, \widehat{y}_{i}\) .
Unlike errors, there exist certain dependencies among the residuals. One
dependency is due to the algebraic fact that the average residual is zero. Fur-
ther, there must be at least \(k+2\) observations for there to be variation in the
fit of the plane. If we have only \(k+1\) observations, we can fit a plane to the
 data perfectly, resulting in no variation in the fit. For example, if \(k=1\) , because
two observations determine a line, then at least three observations are required
to observe any deviation from the line. Because of these dependencies, we have
only \(n-( k+1 )\) free, or unrestricted, residuals to estimate the variability about
 the regression plane.
 The positive square root of \(s^{2}\) is our estimator of \(\sigma\) . Using residuals, it can be
expressed as
\[
s \,=\sqrt{\frac{1} {n \,-\, ( k \,+1 )} \sum_{i=1}^{n} e_{i}^{2}} . \tag{3.8}
\]
Because it is based on residuals, we refer to \(s\) as the residual standard deviation.
 The quantity \(s\) is a measure of our " typicalerror." For this reason, \(s\) is also called
 the standard error of the estimate.
The Coefficient of Determination: \(R^{2}\)
To summarize the goodness of fit of the model, as in Chapter \(2\) we partition the
variability into pieces that are " plained" and "" un\&plained" by the regression
\clearpage
3.3 Estimation and Goodness of Fit 83
 fit. Algebraically, the calculations for regression using many variables are similar
 to the case of using only one variable. Unfortunately, when dealing with many
variables, we lose the easy graphical interpretation such as in Figure 2.4.
Begin with the total sum of squared deviations, Total \(S S \,=\, \sum_{i=1}^{n} \left( y_{i} \,-\, \overline{{y}} \right)^{2}\)
 as our measure of the total variation in the dataset. As in equation (2.1), we may
then interpret the equation
\[
\begin{array} {c} {\underbrace{y_{i}-\overline{{y}}}_{\mathrm{t o t a l}} \ =\underbrace{y_{i}-\widehat{y}_{i}}_{\mathrm{d e v i a t i o n}} \ +\underbrace{\widehat{y}_{i}-\overline{{y}}}_{\mathrm{d e v i a t i o n}}} \\ \end{array}
\]
 as the " deviation without knowledge of the explanatory variables equals the
 deviation not explained by the explanatory variables plus deviation explained by
 the explanatory variables." Squaring each side and summing over all observations
yields
:
\[
\mathit{T o t a l} \; \mathit{S S}=\mathit{E r r o r} \; \mathit{S S}+\mathit{R e g r e s s i o n} \; \mathit{S S} ,
\]
 where Error \(S S \,=\, \sum_{i \,=\, 1}^{n} \, ( y_{i} \,-\, \widehat{y}_{i} )^{2}\) and Regression \(S S \,=\, \sum_{i=1}^{n} \left( \widehat{y_{i}} \,-\, \overline{{y}} \right)^{2}\) As in
Section 2.3 for the one explanatory variable case, the sum of the cross-product
terms turns out to be zero.
 A statistic that summarizes this relationship is the coeffcient of determination,
\[
R^{2}=\cfrac{R e g r e s s i o n \; S S} {T o t a l \; S S} .
\]
We interpret \(R^{2}\) to be the proportion of variability explained by the regression
function.
If the model is a desirable one for the data, one would expect a strong rela-
tionship between the observed responses and those " pected" under the model,
 the fitted values. An interesting algebraic fact is the following. If we square the
correlation coefficient between the responses and the fitted values, we get the
coefficient of determination; that is,
\[
R^{2}=\left[ r \left( y , \, \widehat{y} \right) \right]^{2} .
\]
As a result, \(R\) , the positive square root of \(R^{2}\) is called the multiple correla-
tion coeffcient . It can be interpreted as the correlation between the response
and the best linear combination of the explanatory variables, the fitted values.
(This relationship is developed using matrix algebra in the technical supplement
Section 5.10.1.)
 The variability decomposition is also summarized using the analysis of vari-
ance, or ANOVA, table, as follows:
ANOVA Table
Source Sum of Squares \(d f\) Mean Square
Regression Regression \(\mathit{S S}\) \(k\) Regression Ms
Error Error Ss \(n-( k+1 )\) \(M S E\)
Total Total \(\mathit{S S}\) \(n-1\)
\clearpage
84 Multiple Linear Regression -- I
 Table 3.3 Term Life
ANOVA Table Source Sum of Squares \(d f\) Mean Square
Regression 328.47 3 109.49
Error 630.43 271 2.326
Total 958.90 274
The mean square column figures are defined to be the sum of squares figures
divided by their respective degrees of freedom. The error degrees of freedom
denotes the number of unrestricted residuals. It is this number that we use in our
definition of the " average," or mean, square error. That is, we define
\[
M S E=E r r o r \ M S=\cfrac{E r r o r \ S S} {n-( k+1 )}=s^{2} .
\]
 Similarly, the regression degrees of freedom is the number of explanatory vari-
ables. This yields
\[
R e g r e s s i o n \, \, M S=\cfrac{R e g r e s s i o n \, \, S S} {k} .
\]
When discussing the coefficient of determination, it can be established that when-
 ever an explanatory variable is added to the model, \(R^{2}\) never decreases. This is
 true whether or not the additional variable is useful. We would like a measure of
 fit that decreases when useless variables are entered into the model as explanatory
variables. To circumvent this anomaly, a widely used statistic is the coeffcient of
 determination adjusted for degrees of freedom, defined by
\[
R_{a}^{2}=1-\frac{( E r r o r \ S S ) / [ n-( k+1 ) ]} {( T o t a l \ S S ) / ( n-1 )}=1-\frac{s^{2}} {s_{y}^{2}} . \tag{3.9}
\]
To interpret this statistic, note that \(s_{y}^{2}\) does not depend on the model or on the
model variables. Thus, \(s^{2}\) and \(R_{a}^{2}\) are equivalent measures of model fit. As the
model fit improves, then \(R_{a}^{2}\) becomes larger and \(s^{2}\) becomes smaller, and vice
versa. Put another way, choosing a model with the smallest \(s^{2}\) is equivalent to
choosing a model with the largest \(R_{a}^{2}\) .
Example: Term Life Insurance, Continued. To illustrate, Table 3.3 displays the
 summary statistics for the regression of LNFACE on EDUCATION, NUMHH,
and LNINCOME. From the degrees-of-freedom column, we remind ourselves
 that there are three explanatory variables and 275 observations. As measures of
model fit, the coefficient of determination is \(R^{2}=3 4 . 3 7 6 ~ ( \mathrm{=} 3 2 8 . 4 7 7 9 5 8 . 9 0 )\) and
 the residual standard deviation is \(s \,=\, 1 . 5 2 5 \, (=\sqrt{2 . 3 2 6} )\) If we were to attempt to
 estimate the logarithmic face amount without knowledge of the explanatory vari-
 ables EDUCATION, NUMHH, and LNINCOME, then the size of the typical error 
would be \(s_{y}=1 . 8 7 1 \: (=\sqrt{9 5 8 . 9 0 / 2 7 4} )\) . Thus, by taking advantage of our knowl-
 edge of the explanatory variables, we have been able to reduce the size of the typi-
 cal error. The measure of model fit that compares these two estimates of variability
 is the adjusted coefficient of determination, \(R_{a}^{2}=1 \,-\, 2 . 3 2 6 / 1 . 8 7 1^{2}=3 3 . 6 9 6 .\)
\clearpage
3.4 Statistical Inference for a Single Coeffcient 85
Table 3.4
Variable Coefficient \(t\) -Statistic Regression
Intercept 9.904 12.928 Coefficients from a
 Model of Female
Logarithmic number of persons per physician -0.473 3.212 Advantage
Fertility -0.444 -3.477
Percentage of Hindus and Buddhists \(- 0 . 0 1 8\) -3.196
 Soviet Union dummy 4.922 7.235
Source: Lemaire (2002)
Example: Why Do Females Live Longer Than Males? In an article with this
 title, Lemaire (2002) examined what he called the " femaleadvantage," the differ-
 ence in life expectancy between women and men. Life expectancies are of interest
because they are widely used measures of a nation's health. Lemaire examined
 data from \(n=1 6 9\) countries and found that the average female advantage was
 4.51 years worldwide. He sought to explain this difference based on 45 behavior-
 ial measures, variables that capture a nation's degree of economic modernization;
social, cultural, and religious mores; geographic position; and quality of health
 care available.
 After a detailed analysis, Lemaire reports coefficients from a regression model
that appear in Table 3.4. This regression model explains \(R^{2}=6 1^{\mathcal{V}_{o}}\) of the vari-
ability. It is a parsimonious model consisting of only \(k=4\) of the original 45
variables.
All variables were strongly statistically significant. The number of persons
per physician was also correlated with other variables that capture a country's
degree of economic modernization, such as urbanization, number of cars, and
 percentage working in agriculture. Fertility, the number of births per woman, was
 highly correlated with education variables in the study, including female illiteracy
and female school enrollment. The percentage of Hindus and Buddhists is a
 social, cultural, and religious variable. The Soviet Union dummy is a geographic
variable -- it characterizes Eastern European countries that formerly belonged
to the Soviet Union. Because of the high degree of collinearity among the \(4 5\)
 candidate variables, other analysts could easily pick an alternative set of variables.
Nonetheless, Lemaire's important point was that this simple model explains
roughly \(6 1 7_{o}\) of the variability from only behaviorial variables, unrelated to
 biological sex differences.
3.4 Statistical Inference for a Single Coeffcient
3.4.1 The t-Test
In many applications, a single variable is of primary interest, and other variables
 are included in the regression to control for additional sources of variability. To
illustrate, a sales agent might be interested in the effect that income has on the
\clearpage
86 Multiple Linear Regression -- I
quantity of insurance demanded. In a regression analysis, we could also include
other explanatory variables such as an individual's sex, occupation, age, size of
household, education level, and so on. By including these additional explanatory
variables, we hope to gain a better understanding of the relationship between
 income and insurance demand. To reach sensible conclusions, we will need some
rules to decide whether a variable is important.
We respond to the question, Is \(x_{j}\) important? by investigating whether the
corresponding slope parameter, \(\beta_{j}\) equals zero. The question is whether \(\beta_{j}\) is
 zero can be restated in the hypothesis testing framework as  Is \(H_{0} : \beta_{j}=0\)
valid?
We examine the proximity of \(b_{j}\) to zero to determine whether \(\beta_{j}\) is zero.
Begnae \(b_{j}\) depend on the units of \(y\) and \(x_{j}\) , we need to standardize \(\sigma^{2}\) times
\(( j+1 )\) this quantity. In Property \(2\) and equation (3.6), we saw that Var \(\sigma^{2}\) by the estimator \(b_{j}\) is \(s^{2}\) and
st diagonal element of \(\left( \mathbf{X}^{\prime} \mathbf{X} \right)^{-1}\) Replacing
taking square roots, we have the following:
Defnition . The standard error of \(b_{j}\) can be expressed as
\[
s e ( b_{j} )=s \sqrt{( j+1 ) s t \ d i a g o n a l \ e l e m e n t \ o f \left( \mathbf{X^{\prime} X} \right)^{-1}} .
\]
Recall that a standard error is an estimated standard deviation. To test \(H_{0} : \beta_{j}=0 ,\)
\(t\) -ratio, \(t ( b_{j} )=b_{j} / s e ( b_{j} )\) . We interpret \(t ( b_{j} )\) to be the number
Interpret t(bj) to be of standard errors that \(b_{j}\) iayfrmpre aiy
the number of because the sampling distribution of \(t ( b_{j} )\) can be shown to be the \(t\)
 is away from zero. with \(d f=n-( k+1 )\) degrees of freedom, under the null hypothesis with the
 linear regression model assumptions F1-- F5.This enables us to construct tests of
 the null hypothesis such as the following procedure:
Procedure. The t-test for a Regression Coefficient (beta).
: The null hypothesis is \(H_{0} : \beta_{j}=0 .\)
: The alternative hypothesis is \(H_{a} : \beta_{j} \neq0 .\)
Establish a significance level  (typically but not necessarily 5\%).
: Construct the statistic, \(t ( b_{j} )=b_{j} / s e ( b_{j} )\) .
- Procedure: Reject the null hypothesis in favor of the alternative if \(| t ( b_{j} ) |\)
exceeds a \(t\) -value. Here, this \(t\) -value is the \(( 1-\alpha/ 2 )\) th percentile from
the \(t\) -distribution using \(d f=n \,-\, ( k \,+\, 1 )\) degrees of freedom, denoted as
\(t_{n-( k+1 ) , 1-\alpha/ 2} .\)
In many applications, the sample size will be large enough so that we may
Rule of thumb: approximate the \(t\) -value by the corresponding percentile from the standard normal
inerpret a variable tocurve. At the \(5 7_{o}\) level of significance, this percentile is 1.96. Thus, as a rule of
be important if its
t-ratio exceeds two in thumb, we can interpret a variable to be important if its \(t\) -ratio exceeds two in
 absolute value. absolute value.
\clearpage
3.4 Statistical Inference for a Single Coeffcient 87
Table 3.5
Alternative Procedure: Reject \(H_{0}\) in Decision-Making
I \(H_{a}\) favor of \(H_{a}\) Procedures for
)
\(\beta_{j} < d\) \(\vert t \mathrm{-r a t i o} \vert\, > \, t_{n-( k+1 ) , 1-\alpha/ 2}\) Testing \(H_{0} : \beta_{j}=d\)
\(\beta_{j} > d\) \(t \mathrm{-r a t i o} \, > \, t_{n-\left( k+1 \right) , 1-\alpha}\)
\(\beta_{j} \neq d\) \(t \mathrm{-r a t i o} <-{t}_{n-\left( k+1 \right) , 1-\alpha}\)
Notes: The significance level is \(\alpha\) .Here, \(t_{n-( k+1 ) , 1-\alpha}\)
is the \(( 1 \mathrm{-} \alpha)\) th percentile from the \(t\) -distribution
using \(d f=n \,-\, ( k \,+\, 1 )\) degrees of freedom. The test
statistic is \(t \mathrm{-r a t i o}=( b_{j} \,-d ) / s e ( b_{j} ) .\)
Table 3.6
Alternative Probability Values for
Hypothesis ( \(( H_{a} )\) \(\beta_{j} > d\)
\[
\beta_{j} < d
\]
\(\beta_{j} \neq d\) Testing \(H_{0} : \beta_{j}=d\)
\(p\) -Value \(\operatorname* {P r} ( t_{n-( k+1 )} > t \mathrm{-r a t i o} )\) \(\operatorname* {P r} ( t_{n-( k+1 )} < t \mathrm{-r a t i o} )\) \(\operatorname* {P r} ( | t_{n-( k+1 )} | > | t |\) -ratio()
Notes: Here, \(t_{n-( k+1 )}\) is a \(t\) -distributed random variable with \(d f=n \,-\, ( k \,+\, 1 )\) degrees of
freedom. The test statistic is \(t \mathrm{-r a t i o}=( b_{j} \,-d ) / s e ( b_{j} )\)
Although it is the most common, testing \(H_{0} : \beta_{j}=0\) versus \(H_{a} : \beta_{j} \neq0\) is just
 one of many hypothesis tests that can be performed. Table 3.5 outlines alternative
 decision-making procedures. These procedures are for testing \(H_{0} : \beta_{j}=d\) . Here,
\(d\) is a user-prescribed value that may be equal to zero or any other known value.
Alternatively, one can construct \(p\) -values and compare them to given signif-
icant levels. The \(p\) -value allows the report reader to understand the strength of
the deviation from the null hypothesis. Table 3.6 summarizes the procedure for
calculating \(p\) -values.
Example: Term Life Insurance, Continued. \(\mathrm{A}\) useful convention when report-
 ing the results of a statistical analysis is to place the standard error of a statistic in
 parentheses below that statistic. Thus, for example, in our regression of LNFACE
 on EDUCATION, NUMHH, and LNINCOME, the estimated regression equation
is
\[
\begin{array} {l} {\widehat{\mathrm{L N F A C E}}} \\ {\mathrm{s t a n d a r d ~ e r r o r} \hfill\quad\quad( 0 . 8 4 6 ) \hfill\quad\quad( 0 . 0 3 9 ) \hfill\quad\quad\quad\quad( 0 . 0 6 3 ) \hfill\quad\quad\quad( 0 . 0 7 8 )} \\ \end{array}
\]
To illustrate the calculation of the standard errors, first note that, from
Table 3.3, the residual standard deviation is \(s \,=\, 1 . 5 2 5\) - Using a statistical pack-
 age, we have
\[
\left( \mathbf{X}^{\prime} \mathbf{X} \right)^{-1}=\left( \begin{matrix} {\; \; \, 0 . 3 0 7 9 7 5} & {-0 . 0 0 4 6 3 3} & {-0 . 0 0 2 1 3 1} & {-0 . 0 2 0 6 9 7} \\ {-0 . 0 0 4 6 3 3} & {0 . 0 0 0 6 4 8} & {0 . 0 0 0 1 4 3} & {-0 . 0 0 0 4 6 7} \\ {-0 . 0 0 2 1 3 1} & {0 . 0 0 0 1 4 3} & {0 . 0 0 1 7 2 4} & {-0 . 0 0 0 4 5 3} \\ {-0 . 0 2 0 6 9 7} & {-0 . 0 0 0 4 6 7} & {-0 . 0 0 0 4 5 3} & {\; \; \, 0 . 0 0 2 5 8 5} \\ \end{matrix} \right) .
\]

\clearpage
88 Multiple Linear Regression -- I
To illustrate, we can compute \(s e ( b_{3} )=s \times\sqrt{0 . 0 0 2 5 8 5}=0 . 0 7 8 ,\) as earlier. Cal-
culation of the standard errors, as well as the corresponding \(t\) -statistics, is part of
 the standard output from statistical software and need not be computed by users.
 Our purpose here is to illustrate the ideas underlying the routine calculations.
With this information, we can immediately compute \(t\) -ratios to check to
see whether a coefficient associated with an individual variable is significantly
different from zero. For example, the \(t\) -ratio for the LNINCOME variable is
\(t ( b_{3} )=0 . 4 9 4 / 0 . 0 7 8=6 . 3 .\) The interpretation is that \(b_{3}\) is more than four stan-
dard errors above zero, and thus LNINCOME is an important variable in the
model. More formally, we may be interested in testing the null hypothesis that
\(H_{0} : \beta_{3}=0\) versus \(H_{0} : \beta_{3} \neq0 .\) At a 5\% level of significance, the \(t\) -value is
1.96, because \(d f=2 7 5-\left( 1+3 \right)=2 7 1\) . We thus reject the null in favor of the
alternative hypothesis, that logarithmic income (LNINCOME) is important in
 determining the logarithmic face amount.
3.4.2 Confidence Intervals
Confdence intervals for parameters provide another device for describing the
strength of the contribution of the \(j\) th explanatory variable. The statistic \(b_{j}\) is a
 point estimate of the parameter \(\beta_{j}\) To provide a range of reliability, we use the
confidence interval
\[
b_{j} \pm t_{n-( k+1 ) , 1-\alpha/ 2} s e ( b_{j} ) . \tag{3.10}
\]
Here, the \(t\) -value \(t_{n-( k+1 ) , 1-\alpha/ 2}\) is a percentile from the \(t\) -distribution with \(d f=\)
\(n-( k+1 )\) degrees of freedom. We use the same \(t\) -value as in the two-sided
hypothesis test. Indeed, there is a duality between the confidence interval and
the two-sided hypothesis test. For example, it is not hard to check that if a
 hypothesized value falls outside the confidence interval, then \(H_{0}\) will be rejected
 in favor of \(H_{a}\) Further, knowledge of the \(p\) -value, point estimate, and standard
error can be used to determine a confidence interval.
 3.4.3 Added Variable Plots
To represent multivariate data graphically, we have seen that a scatterplot matrix
is a useful device. However, the major shortcoming of the scatterplot matrix is
 that it captures relationships only between pairs of variables. When the data can
be summarized using a regression model, a graphical device that does not have
 this shortcoming is an added variable plot. The added variable plot is also called
a partial regression plot because, as we will see, it is constructed in terms of
residuals from certain regression fits. We will also see that the added variable plot
can be summarized in terms of a partial correlation coefficient, thus providing a
 link between correlation and regression. To introduce these ideas, we work in the
context of the following example:
\clearpage
3.4 Statistical Inference for a Single Coeffcient 89
Table 3.7 Summary
Standard Statistics for Each
Variable Mean Median Deviation Minimum Maximum Variable for 37
Refrigerators
ECOST 70.51 68.00 9.14 60.00 94.00
RSIZE 13.400 13.200 0.600 12.600 14.700
FSIZE 5.184 5.100 0.938 4.100 7.400
 SHELVES 2.514 2.000 1.121 1.000 5.000
FEATURES 3.459 3.000 2.512 1.000 12.000
PRICE 626.4 590.0 139.8 460.0 1200.0
Source: Consumer Reports, July 1992. " Refrigerators:A Comprehensive Guide to the Big
White Box."
Table 3.8 Matrix
ECOST RSIZE FSIZE SHELVES FEATURES of Correlation
Coefficients
RSIZE 0.333
FSIZE 0.855 0.235
SHELVES 0.188 0.363 0.251
FEATURES 0.334 0.096 \(0 . 4 3 9\) 0.160
PRICE 0.522 0.024 0.720 0.400 0.697
 EMPIRICAL
Example: Refrigerator Prices. What characteristics of a refrigerator are impor- Filename is
tant in determining its price (PRICE)? We consider here several characteristics "Refrigerator"
of a refrigerator, including the size of the refrigerator in cubic feet (RSIZE), the
size of the freezer compartment in cubic feet (FSIZE), the average amount of
money spent per year to operate the refrigerator (ECOST, for " enegy cost" ),
the number of shelves in the refrigerator and freezer doors (SHELVES), and
 the number of features (FEATURES). The features variable includes shelves for
cans, see-through crispers, ice makers, egg racks, and so on.
Both consumers and manufacturers are interested in models of refrigerator
prices. Other things equal, consumers generally prefer larger refrigerators with
lower energy costs that have more features. Because of forces of supply and
 demand, we would expect consumers to pay more for such refrigerators. A larger
refrigerator with lower energy costs that has more features at the similar price is
considered a bargain to the consumer. How much extra would the consumer be
willing to pay for the additional space? A model of prices for refrigerators on the
market provides some insight to this question.
 To this end, we analyze data from \(n=3 7\) refrigerators. Table 3.7 provides the
 basic summary statistics for the response variable PRICE and the five explanatory
variables. From this table, we see that the average refrigerator price is \(\overline{{y}}=\)
\$626.40, with standard deviation \(s_{y}=\S1 3 9 . 8 0\) Similarly, the average annual
amount to operate a refrigerator, or average ECOST, is \$70.51.
 To analyze relationships among pairs of variables, Table 3.8 provides a matrix
of correlation coefficients. From the table, we see that there are strong linear
relationships between PRICE and each of freezer space (FSIZE) and the number
\clearpage
90 Multiple Linear Regression -- I
Table 3.9 Fitted
Refrigerator Price Coefficient Standard
 Model Estimate Error \(t\) -Ratio
Intercept 798 271.4 2.9
ECOST 6.96 2.275 \(- 3 . 1\)
RSIZE 76.5 \(1 9 . 4 4\) 3.9
FSIZE 137 23.76 5.8
SHELVES 37.9 9.886 3.8
FEATURES 23.8 4.512 5.3
 of FEATURES. Surprisingly, there is also a strong positive correlation between
PRICE and ECOST. Recall that ECOST is the energy cost; one might expect that 
 higher-priced refrigerators should enjoy lower energy costs.
 A regression model was fit to the data. The fitted regression equation appears
 in Table 3.9, with \(s \,=6 0 . 6 5\) and \(R^{2}=8 3 . 8 \mathcal{V}_{o}\) .
From Table 3.9, the explanatory variables seem to be useful predictors of
refrigerator prices. Together, the variables account for 83.8\% of the variability. To
 understand prices, the typical error has dropped from \(s_{y} \,=\, \S1 3 9 . 8 0 \, \mathrm{t o} \, s \,=\, \S6 0 . 6 5 .\)
The \(t\) -ratios for each of the explanatory variables exceeds 2 in absolute value,
 indicating that each variable is important on an individual basis.
What is surprising about the regression fit is the negative coefficient associated
with energy cost. Remember, we can interpret \(b_{E C O S T}=-6 . 9 6\) to mean that, for
each dollar increase in ECOST, we expect the PRICE to decrease by \$6.96. This
 negative relationship conforms to our economic intuition. However, it is surprising
that the same dataset has shown us that there is a positive relationship between
PRICE and ECOST. This seeming anomaly is because correlation measures
relationships only between pairs of variables, though the regression fit can account
for several variables simultaneously. To provide more insight into this seeming
 anomaly, we now introduce the added variable plot.
 Producing an Added Variable Plot
 The added variable plot provides additional links between the regression method-
 ology and more fundamental tools such as scatter plots and correlations. We work
 in the context of the refrigerator price example to demonstrate the construction
of this plot.
Procedure for producing an added variable plot.
(i) Run a regression of PRICE on RSIZE, FSIZE, SHELVES, and FEA-
TURES, omitting ECOST. Compute the residuals from this regression,
 which we label \(e_{1}\) .
(ii) Run a regression of ECOST on RSIZE, FSIZE, SHELVES, and
FEATURES. Compute the residuals from this regression, which we
label \(e_{2}\) .
(iii) Plot \(e_{1}\) versus \(e_{2}\) . This is the added variable plot of PRICE versus
ECOST, controlling for the effects of the RSIZE, FSIZE, SHELVES, and
 FEATURES. This plot appears in Figure 3.4.
\clearpage
3.4 Statistical Inference for a Single Coeffcient 91
150 0 \(\circ\) Figure 3.4 An added
variable plot. The
100 \(\varpi\) 0 residuals from the
regression of PRICE
0 \(\circ\surd\omicron\) \(\circ\) : 0 0 variables, omitting
\(\circ\) on the explanatory
50 \(\circ^{\circ}\) : ECOST, are on the
0 8 \(\circ_{\mathrm{O}}\) vertical axis. On the
0 horizontal axis are the
100 \(\circ\ss_{\omicron}\) 0 regression fit of
residuals from the
0 ECOST on the other
5 \(0\) 5 10 explanatory variables.
The correlation
\(\mathbf{e}_{1}\) coefficient is \(- 0 . 4 8\) .
The error \(\varepsilon\) can be interpreted as the natural variation in a sample. In many
situations, this natural variation is small compared to the patterns evident in
the nonrandom regression component. Thus, it is useful to think of the error,
\(\varepsilon_{i} \,=\, y_{i} \,-( \beta_{0}+\beta_{1} x_{i \, 1}+\cdots+\beta_{k} x_{i \, k} ) .\) as the response after controlling for the
effects of the explanatory variables. In Section 3.3, we saw that a random error
can be approximated by a residual, \(e_{i} \,=\, y_{i} \,-\, ( b_{0} \,+\, b_{1} x_{i \, 1} \,+\, \cdots\,+\, b_{k} \, x_{i \, k} ) .\) Thus,
 in the same way, we may think of a residual as the response after " controlling
for" the effects of the explanatory variables.
 With this in mind, we can interpret the vertical axis of Figure 3 .4 as the refriger-
 ator PRICE controlled for effects of RSIZE, FSIZE, SHELVES, and FEATURES.
 Similarly, we can interpret the horizontal axis as the ECOST controlled for effects
 of RSIZE, FSIZE, SHELVES, and FEATURES. The plot then provides a graph-
 ical representation of the relation between PRICE and ECOST, after controlling
for the other explanatory variables. For comparison, a scatter plot of PRICE and
 ECOST (not shown here) does not control for other explanatory variables. Thus,
 it is possible that the positive relationship between PRICE and ECOST is due not
 to a causal relationship but rather to one or more additional variables that cause
 both variables to be large.
For example, from Table 3.7, we see that the freezer size (FSIZE) is positively
correlated with both ECOST and PRICE. It certainly seems reasonable that
 increasing the size of a freezer would cause both the energy cost and the price to
 increase. Rather, the positive correlation may be because large values of FSIZE
 mean large values of both ECOST and PRICE.
Variables left out of a regression are called omitted variables. This omission
could cause a serious problem in a regression model fit; regression coefficients
could be not only strongly significant when they should not be but also of the
 incorrect sign. Selecting the proper set of variables to be included in the regression
model is an important task; it is the subject of Chapters \(5\) and \(6\) .
 3.4.4 Partial Correlation Coefficients
 As we saw in Chapter \(2\) a correlation statistic is a useful quantity for summarizing
 plots. The correlation for the added variable plot is called a partial correlation
coeffcient . It is defined to be the correlation between the residuals \(e_{1}\) and \(e_{2}\) and is
\clearpage
92 Multiple Linear Regression -- I
 denoted by \(r ( y , \, x_{j} | x_{1} , \, \ldots, \, x_{j-1} , \, x_{j+1} , \, \ldots, \, x_{k} )\) Because it summarizes an added
variable plot, we may interpret \(r ( y , \, x_{j} | x_{1} , \, \ldots, \, x_{j-1} , \, x_{j+1} , \, \ldots, \, x_{k} ) )\) to be the
 correlation between \(y\) and \(x_{j}\) , in the presence of the other explanatory variables.
To illustrate, the correlation between PRICE and ECOST in the presence of the
 other explanatory variables is \(- \, 0 . 4 8\) .
 The partial correlation coefficient can also be calculated using
\[
r ( y , \, x_{j} | x_{1} , \, \dots, \, x_{j-1} , \, x_{j+1} , \, \dots, \, x_{k} )=\cfrac{t ( b_{j} )} {\sqrt{t ( b_{j} )^{2}+n-( k+1 )}} . \tag{3.11}
\]
Here, \(t ( b_{j} )\) is the \(t\) -ratio for \(b_{j}\) from a regression of \(y\) on \(x_{1} , \dots, \, x_{k}\) (including the
variable \(x_{j}\) ). An important aspect of equation (3.11) is that it allows us to calculate
partial correlation coefficients running only one regression. For example, from
 Table 3.9, the partial correlation between PRICE and ECOST in the presence of
 the other explanatory variables is \((-3 . 1 ) / \sqrt{(-3 . 1 )^{2} \,+\, 3 7 \,-\, ( 5 \,+\, 1 )} \stackrel{\_} \approx-0 . 4 8\) :
Calculation of partial correlation coefficients is quicker when using the rela-
tionship with the \(t\) -ratio, but may fail to detect nonlinear relationships. The
 information in Table 3.9 allows us to calculate all five partial correlation coef-
ficients in the refrigerator price example after running only one regression. The
 three-step procedure for producing added variable plots requires ten regressions,
two for each of the five explanatory variables. Of course, by producing added
variable plots, we can detect nonlinear relationships that are missed by correlation
coefficients.
Partial correlation coefficients provide another interpretation for \(t\) -ratios.
Equation (3.11) shows how to calculate a correlation statistic from a \(t\) -ratio, thus
 providing another link between correlation and regression analysis. Moreover,
from equation (3.11), we see that the larger is the \(t\) -ratio, the larger is the partial
 correlation coefficient. That is, a large \(t\) -ratio means that there is a large correlation
 between the response and the explanatory variable, controlling for other explana-
tory variables. This provides a partial response to the question that is regularly
asked by consumers of regression analyses: Which variable is most important?
 3.5 Some Special Explanatory Variables
The linear regression model is the basis of a rich family of models. This section
The linear regression provides several examples to illustrate the richness of this family. These examples
 function is linear in demonstrate the use of (i) binary variables, (ii) transformation of explanatory
the parameters but variables, and (iii) interaction terms. This section also serves to underscore the
 may be a highly meaning of the adjective linear in the phrase " linearregression" ;the regression
 nonlinear function of function is linear in the parameters but may be a highly nonlinear function of the
the explanatory
variables. explanatory variables.
 3.5.1 Binary Variables
 Categorical variables provide a numerical label for measurements of observations
that fall in distinct groups, or categories. Because of the grouping, categorical
\clearpage
 3.5 Some Special Explanatory Variables 93
LNFACE Figure 3.5 Letter
16 plot of LNFACE
LNFACE = 5.09 + 0.634 LNINCOME versus LNINCOME,
14 \(\circ\) \(\mathcal{s}\) \(\circ\) Wwih tre lener ode o
" S"for single and \(\varsigma\)
for other. The fitted
12 regression lines have
been superimposed.
10 \(\delta\) \(\circ\) The lower line is for
8 s S S \_NFACE. single and the upper
= 4.29 + 0.634 LNINCOme line is for other.
6 8 10 12
LNINCOME
variables are discrete and generally take on a finite number of values. We begin
our discussion with a categorical variable that can take on one of only two
values, a binary variable. Further discussion of categorical variables is the topic
 of Chapter \(4\) .
Example: Term Life Insurance, Continued. We now consider the marital status
of the survey respondent. In the Survey of Consumer Finances, respondents can
 select among several options describing their marital status including " married,
" living with a partner," " divorced," and so on. Marital status is not measured
continuously but rather takes on values that fall into distinct groups. In this Binary explanatory
chapter, we group survey respondents according to whether they are single, variables are also
 known as indicator
defined to include those who are separated, divorced, widowed, never married, and dummy variables.
 and are not married nor living with a partner. Chapter \(4\) will present a more
 complete analysis of marital status by including additional categories.
 The binary variable SINGLE is defined as \(1\) if the survey respondent is single
and \(\; \; \; \; \; \; \; \; \; \; \; \; \; \; \; \; \; \; \; \; \; \; \; \; \; \; \; \; \; \; \; \; \; \; \; \; \; \; \; \; \; \; \; \; \; \; \; \; \; \; \; \; \; \; \; \; \; \; \; \; \; \; \; \; \; \; \; \; \; \; \; \; \; \; \; \; \; \; \; \; \; \; \; \; \; \; \; \; \; \; \; \; \; \; \; \; \; \; \; \; \; \; \; \; \; \; \; \; \; \; \; \; \; \; \; \; \; \; \; \; \; \; \; \; \; \; \; \; \; \; \; \; \; \; \; \; \; \; \; \; \; \; \; \; 0 \; \; \; \; \; 0\) otherwise. The variable SINGLE is also known as an indicator variable
because it indicates whether the respondent is single. Another name for this
 important type of variable is a dummy variable. We could use \(\; \; \; \; \; \; \; \; \; \; \; \; \; \; \; \; \; \; \; \; \; \; \; \; \; \; \; \; \; \; \; \; \; \; \; \; \; \; \; \; \; \; \; \; \; \; \; \; \; \; \; \; \; \; \; \; \; \; \; \; \; \; \; \; \; \; \; \; \; \; \; \; \; \; \; \; \; \; \; \; \; \; \; \; \; \; \; \; \; \; \; \; \; \; \; \; \; \; \; \; \; \; \; \; \; \; \; \; \; \; \; \; \; \; \; \; \; \; \; \; \; \; \; \; \; \; \; \; \; \; \; \; \; \; \; \; \; \; \; \; \; \; \; \; 0 \; \; \; \; \; 0\) and 100, or 20
 and 36, or any other distinct values. However, \(\; \; \; \; \; \; \; \; \; \; \; \; \; \; \; \; \; \; \; \; \; \; \; \; \; \; \; \; \; \; \; \; \; \; \; \; \; \; \; \; \; \; \; \; \; \; \; \; \; \; \; \; \; \; \; \; \; \; \; \; \; \; \; \; \; \; \; \; \; \; \; \; \; \; \; \; \; \; \; \; \; \; \; \; \; \; \; \; \; \; \; \; \; \; \; \; \; \; \; \; \; \; \; \; \; \; \; \; \; \; \; \; \; \; \; \; \; \; \; \; \; \; \; \; \; \; \; \; \; \; \; \; \; \; \; \; \; \; \; \; \; 0 \; \; \; \; \; 0\) and \(1\) are convenient for the
 interpretation of the parameter values, discussed subsequently. To streamline the
 discussion, we now present a model using only LNINCOME and SINGLE as
 explanatory variables.
 For our sample of \(n \,=2 7 5\) households, 57 are single and the other 218 are not.
 To see the relationships among LNFACE, LNINCOME, and SINGLE, Figure 3.5
introduces a letter plot of LNFACE versus LNINCOME, with SINGLE as the
code variable. We can see that Figure 3.5 is a scatter plot of LNFACE versus
LNINCOME, using \(5 0\) randomly selected households from our sample of 275
(for clarity of the graph). However, instead of using the same plotting symbol for
 each observation, we have coded the symbols so that we can easily understand the
behavior of a third variable, SINGLE. In other applications, you may elect to use
 other plotting symbols such as \(\clubsuit, \heartsuit, \spadesuit\) and so on, or to use different colors, to
\clearpage
94 Multiple Linear Regression -- I
encode additional information. For this application, the letter codes " S"for single
and "- o"for other were selected because they remind the reader of the plot of the
nature of the coding scheme. Regardless of the coding scheme, the important
point is that a letter plot is a useful device for graphically portraying three or
more variables in two dimensions. The main restriction is that the additional
information must be categorized, such as with binary variables, to make the
coding scheme work.
Figure 3.5 suggests that LNFACE is lower for those who are single than for
others for a given level of income. Thus, we now consider a regression model,
\[
\qquad\qquad\qquad L N F A C E=\beta_{0} \,+\, \beta_{1} L N I N C O M E \,+\, \beta_{2} S I N G L E+\, \varepsilon.
\]
 The regression function can
be written as
\[
\textrm{E} y=\begin{cases} {\beta_{0}+\beta_{1} \mathrm{L N I N C O M E}} & {\mathrm{f o r ~ o t h e r ~ r e s p o n d e n t s}} \\ {\beta_{0}+\beta_{2}+\beta_{1} \mathrm{L N I N C O M E}} & {\mathrm{f o r ~ s i n g l e ~ r e s p o n d e n t s}} \\ \end{cases} .
\]
The interpretation of the model coefficients differs from the continuous vari-
able case. For continuous variables such as LNINCOME, we interpret \(\beta_{1}\) as
the expected change in \(y\) per unit change of logarithmic income, holding other
variables fixed. For binary variables such as SINGLE, we interpret \(\beta_{2}\) as the
expected increase in \(y\) when going from the base level of SINGLE ( \(= 0 )\) to the
 alternative level. Thus, although we have one model for both marital statuses,
we can interpret the model using two regression equations, one for each type of
marital status. By writing a separate equation for each marital status, we have
been able to simplify a complicated multiple regression equation. Sometimes,
 you will find that it is easier to communicate a series of simple relationships than
a single, complex relationship.
 Although the interpretation for binary explanatory variables differs from the
continuous, the ordinary least squares estimation method remains valid. To illus-
trate, the fitted version of the preceding model is
\[
\begin{array} {l} {\widehat{L N F A C E}} \\ {\mathrm{s t a n d a r d ~ e r r o r} \quad( 0 . 8 9 ) \qquad\qquad( 0 . 0 7 8 )} \\ \end{array} \begin{array} {l} {{\mathrm{~}}} \\ {( 0 . 2 4 8 )} \\ \end{array} .
\]
 To interpret \(b_{2}=-0 . 8 0 0\) we say that we expect the logarithmic face to be smaller
 by 0.80 for a survey respondent who is single compared to the other category. This
assumes that other things, such as income, remain unchanged. For a graphical
interpretation, the two fitted regression lines are superimposed in Figure 3.5.
3.5.2 Transforming Explanatory Variables
Regression models have the ability to represent complex, nonlinear relationships
 between the expected response and the explanatory variables. For example, early
regression texts, such as Plackett (1960, chapter 6) devote an entire chapter of
 material to polynomial regression,
\[
\textrm{E} y=\beta_{0}+\beta_{1} x+\beta_{2} x^{2}+\cdots+\beta_{p} x^{p} . \tag{3.12}
\]

\clearpage
 3.5 Some Special Explanatory Variables 95
Here, the idea is that a \(p\) th order polynomial in \(x\) can be used to approximate
 general, unknown nonlinear functions of \(x\) .
 The modern-day treatment of polynomial regression does not require an entire
chapter because the model in equation (3.12) can be expressed as a special
case of the linear regression model. That is, with the regression function in
equation (3.5), E \(y=\beta_{0}+\beta_{1} x_{1}+\beta_{2} x_{2}+\cdots+\beta_{k} x_{k}\) , we can choose \(k=p\) and
\(x_{1}=x , \, x_{2}=x^{2} , \, \dots, \, x_{p}=x^{\, p}\) . Thus, with these choices of explanatory variables,
 we can model a highly nonlinear function of \(x\) .
We are not restricted to powers of \(x\) in our choice of transformations. For
example, the model E \(y=\beta_{0}+\beta_{1} \operatorname{l n} x\) , provides another way to represent a
 gently sloping curve in \(x\) . This model can be written as a special case of the basic
 linear regression model using \(x^{*}=\operatorname{l n} x\) as the transformed version of \(x\) .
Transformations of explanatory variables need not be smooth functions. To
 illustrate, in some applications, it is useful to categorize a continuous explanatory
 variable. For example, suppose that \(x\) represents the number of years of education,
ranging from \(\; \; \; \; \; \; \; \; \; \; \; \; \; \; \; \; \; \; \; \; \; \; \; \; \; \; \; \; \; \; \; \; \; \; \; \; \; \; \; \; \; \; \; \; \; \; \; \; \; \; \; \; \; \; \; \; \; \; \; \; \; \; \; \; \; \; \; \; \; \; \; \; \; \; \; \; \; \; \; \; \; \; \; \; \; \; \; \; \; \; \; \; \; \; \; \; \; \; \; \; \; \; \; \; \; \; \; \; \; \; \; \; \; \; \; \; \; \; \; \; \; \; \; \; \; \; \; \; \; \; \; \; \; \; \; \; \; \; \; \; \; \; \; \; \; \; \; \; 0 \; \; \; \; 0\) to 17. If we are relying on information self-reported by our sample
of senior citizens, there may be a substantial amount of error in the measurement
of \(x\) : We could elect to use a less informative but more reliable transform of
\(x\) such as \(x^{*}\) , a binary variable for finishing 13 years of school (finishing high
school). Formally, we would code \(x^{*}=1\) if \(x \geq1 3\) and \(x^{\, *}=0\) if \(x < 1 3\)
Thus, there are several ways that nonlinear functions of the explanatory vari-
 ables can be used in the regression model. An example of a nonlinear regression
 model is \(y=\beta_{0}+\operatorname{e x p} ( \beta_{1} x )+\varepsilon\) These typically arise in science applications of
regressions where there are fundamental scientific principles guiding the complex
 model development.
3.5.3 Interaction Terms
We have so far discussed how explanatory variables, say, \(x_{1}\) and \(x_{2}\) , affect the
mean response in an additive fashion, that is, \(\textrm{E} y \!=\! \beta_{0}+\! \beta_{1} x_{1}+\beta_{2} x_{2}\) . Here, we
expect \(y\) to increase by \(\beta_{1}\) per unit increase in \(x_{1}\) , with \(x_{2}\) held fixed. What if the
 marginal rate of increase of \(\textrm{E} y\) differs for high values of \(x_{2}\) when compared to
low values of \(x_{2}\) ? One way to represent this is to create an interaction variable
\(x_{3}=x_{1} \times x_{2}\) and consider the model E \(\textsc{l} y=\beta_{0}+\beta_{1} x_{1}+\beta_{2} x_{2}+\beta_{3} x_{3} ,\)
With this model, the change in the expected \(y\) per unit change in \(x_{1}\) now
depends on \(x_{2}\) . Formally, we can assess small changes in the regression function
as
\[
\frac{\partial\textbf{E} y} {\partial x_{1}}=\frac{\partial} {\partial x_{1}} \left( \beta_{0}+\beta_{1} x_{1}+\beta_{2} x_{2}+\beta_{3} x_{1} x_{2} \right)=\beta_{1}+\beta_{3} x_{2} .
\]
In this way, we may allow for more complicated functions of \(x_{1}\) and \(x_{2}\) . Figure 3.6
 illustrates this complex structure. From this figure and the preceding calculations,
 we see that the partial changes of \(\textrm{E} y\) due to movement of \(x_{1}\) depend on the value
of \(x_{2}\) . In this way, we say that the partial changes due to each variable are not
unrelated but rather "" move together."
\clearpage
96 Multiple Linear Regression -- I
Figure 3.6 Plot of
\(\textrm{E} y=\beta_{0}+\beta_{1} x_{1}+\)
\(\beta_{2} x_{2}+\beta_{3} x_{1} x_{2}\) versus
\(x_{1}\) and \(x_{2}\) .
80
 60
40 
12
10.2
8
10
X1
6
15
 More generally, an interaction term is a variable that is created as a nonlinear
function of two or more explanatory variables. These special terms, even though
 permitting us to explore a rich family of nonlinear functions, can be cast as special
cases of the linear regression model. To do this, we simply create the variable
of interest and treat this new term as another explanatory variable. Of course,
not every variable that we create will be useful. In some instances, the created
variable will be so similar to variables already in our model that it will provide
no new information. Fortunately, we can use \(t\) -tests to check whether the new
variable is useful. Further, Chapter \(4\) will introduce a test to decide whether a
 group of variables is useful.
The function that we use to create an interaction variable must be more than
 just a linear combination of other explanatory variables. For example, if we use
\(x_{3}=x_{1}+x_{2}\) , we will not be able to estimate all of the parameters. Chapter 5 will
 introduce some techniques to help avoid situations when one variable is a linear
combination of the others.
To give you some exposure to the wide variety of potential applications of
special explanatory variables, we now present a series of short examples.
Example: Term Life Insurance, Continued. How do we interpret the inter-
action of a binary variable with a continuous variable? To illustrate, consider a
Term Life regression model, \(L N F \! A C E=\beta_{0} \,+\, \beta_{1} \mathrm{I}\) LNINCOME \(+ \ \beta_{2} \mathrm{S I N G L E}+\)
\(\beta_{2}\) LNINCOME*SINGLE \(+ \varepsilon\) . In this model, we have created a third explana-
 tory variable through the interaction of LNINCOME and SINGLE. The regression
function can be written as:
\[
\textsc{e} y=\begin{cases} {\beta_{0}+\beta_{1} \textsc{L N I N C O M E}} & {\mathrm{f o r ~ o t h e r ~ r e s p o n d e n}} \\ {\beta_{0}+\beta_{2}+( \beta_{1}+\beta_{3} ) \textsc{L N I N C O M E}} & {\mathrm{f o r ~ s i n g l e ~ r e s p o n d e n}} \\ \end{cases}
\]
ts
nts
Thus, through this single model with four parameters, we can create two separate
regression lines, one for those single and one for others. Figure 3.7 shows the
 two fitted regression lines for our data.
\clearpage
 3.5 Some Special Explanatory Variables 97
Table 3.10
Variable Variable Squared Twenty-Three
Variable \(( D=0 )\) \(( D=1 )\) \(( D=0 )\) eraction with Regression
Baseline Interaction with Baseline \(D=1 )\) Expense Cost Model
Coefficients from an
 Number of life policies issued \(( x_{1}\) \(- 0 . 4 5 4\) 0.152 0.032 0.007
 Amount of term life insurance 0.112 0.206 0.002 0.005
sold \(( x_{2} )\)
Amount of whole life insurance 0.184 0.173 0.008 0.007
sold \(( x_{3} )\)
Total annuity considerations ( \(x_{4}\) ) 0.098 0.169 -0.003 0.009
Total accident and health -0.171 0.014 0.010 0.002
premiums \(( x_{5}\) )
Intercept 7.726
Price of labor (PL) 0.553
Price of capital (PC) 0.102
Note: \(x_{1}\) through \(x_{5}\) are in logarithmic units.
Source: Segal (2002).
LNFACE Figure 3.7 Letter
16- plot of LNFACE
LNFACE = 5.78 + 0.573 LNINCOME versus LNINCOME,
14 : Wih tre lete ode o"
" S"for single and \(\cdot\epsilon\)
for other. The fitted
12 regression lines have
been superimposed.
10 \(\delta\) \(\circ\) The lower line is for
8 INFACE single and the upper
=1.51+1.185 LNINCOME line is for other.
6 8 10 12
LNINCOME
Example: Life Insurance Company Expenses. In a well-developed life insur-
ance industry, minimizing expenses is critical for a company's competitive posi-
 tion. Segal (2002) analyzed annual accounting data from more than 100 firms for
the period 1995- 1998,inclusive, using a data base from the National Associa-
tion of Insurance Commissioners (NAIC) and other reported information. Segal
 modeled overall company expenses as a function of firm outputs and the price of
inputs. The outputs consist of insurance production, measured by \(x_{1}\) through \(x_{5}\)
described in Table 3.10. Segal also considered the square of each output, as well
as an interaction term with a dummy/binary variable \(\mathit{D}\) that indicates whether
or not the firm uses a branch company to distribute its products. (In a branch
company, field managers are company employees, not independent agents.)
For the price inputs, the price of labor (PL) is defined to be the total cost of
employees and agents divided by their number, in logarithmic units. The price
 of capital (PC) is approximated by the ratio of capital expense to the number of
\clearpage
98 Multiple Linear Regression -- I
employees and agents, also in logarithmic units. The price of materials consists of
 expenses other than labor and capital divided by the number of policies sold and
 terminated during the year. It does not appear directly as an explanatory variable.
Rather, Segal took the dependent variable ( \(y )\) to be total company expenses
 divided by the price of materials, again in logarithmic units.
With these variable definitions, Segal estimated the following regression
function:
1
\[
\begin{aligned} {\textnormal{E} y=\beta_{0}+\sum_{j=1}^{5} \left( \beta_{j} x_{j}+\beta_{j+5} D x_{j}+\beta_{j+1 0} x_{j}^{2}+\beta_{j+1 5} D x_{j}^{2} \right)+\beta_{2 1} P L+\beta_{2 2} P C .} \\ \end{aligned}
\]
 The parameter estimates appear in Table 3.10. For example, the marginal change
in \(\textrm{E} y\) per unit change in \(x_{1}\) is
\[
\frac{\partial\textrm{E} y} {\partial x_{1}}=\beta_{1}+\beta_{6} D+2 \beta_{1 1} x_{1}+2 \beta_{1 6} D x_{1} ,
\]
which is estimated as \(- 0 . 4 5 4 \,+\, 0 . 1 5 2 \, D \,+\, ( 0 . 0 6 4 \,-\, 0 . 0 1 4 \, D ) x_{1}\) . For these data,
the median number of policies issued was \(x_{1}=1 5 , 9 4 4\) : At this value of
\(x_{1}\) , the estimated marginal change is \(- 0 . 4 5 4 \,+\, 0 . 1 5 2 \, D \,+\, ( 0 . 0 6 4 \,-\, 0 . 0 1 4 \, D )\)
\(\operatorname{l n} ( 1 5 9 4 4 )=0 . 1 6 5+0 . 0 1 7 D\) or 0.165 for baseline( \(D=0\) and 0.182 for
branch ( \(D=1\) companies.
These estimates are elasticities, as defined in Section 3.2.2. To interpret these
coefficients further, let COST represent total general company expenses and
 NUMPOL represent the number of life policies issued. Then, for branch ( \(D=1\) )
companies, we have
\[
0 . 1 8 2 \approx\frac{\partial y} {\partial x_{1}}=\frac{\partial\, \operatorname{l n} \, C O S T} {\partial\, \operatorname{l n} \, N U M P O L}=\frac{\frac{\partial\, \, C O S T} {\partial\, \, N U M P O L}} {\frac{C O S T} {N U M P O L}} ,
\]
\(\frac{\partial\ C O S T} {\partial\ N U M P O L} \approx0 . 1 8 2 \frac{C O S T} {N U M P O L}\) The median cost is \$15,992,000, so the marginal
or
cos an values is 0.182 x (15992000/15944) = \$182.55.
 Special Case: Curvilinear Response Functions. We can expand the polynomial
 functions of an explanatory variable to include several explanatory variables. For
example, the expected response, or response function, for a second-order model
with two explanatory variables is
\[
\mathrm{E} y=\beta_{0}+\beta_{1} x_{1}+\beta_{2} x_{2}+\beta_{1 1} x_{1}^{2}+\beta_{2 2} x_{2}^{2}+\beta_{1 2} x_{1} x_{2} .
\]
Figure 3.8 illustrates this response function. Similarly, the response function
for a second-order model with three explanatory variables is
\[
\begin{aligned} {\mathrm{E} y} & {{}=\beta_{0}+\beta_{1} x_{1}+\beta_{2} x_{2}+\beta_{3} x_{3}+\beta_{1 1} x_{1}^{2}+\beta_{2 2} x_{2}^{2}+\beta_{3 3} x_{3}^{2}+\beta_{1 2} x_{1} x_{2}} \\ {} & {{} \quad+\, \beta_{1 3} x_{1} x_{3}+\beta_{2 3} x_{2} x_{3} .} \\ \end{aligned}
\]
When there is more than one explanatory variable, third- and higher-order
 models are rarely used in applications.
\clearpage
3.5 Some Special Explanatory Variables 99
Figure 3.8 Plot of
\(\textrm{E} y=\beta_{0}+\beta_{1} x_{1}+1\)
\(\beta_{2} x_{2}+\beta_{1 1} x_{1}^{2}+\)
200 \(\beta_{2 2} x_{2}^{2}+\beta_{1 2} x_{1} x_{2}\) and \(x_{2}\)
versus \(x_{1}\)
150
100
50
0 12
5 \(1 0 \ddag\)
+10 8
6
15
Figure 3.9 The
marginal change in
\(\textrm{E} y\) is lower below
\$97,500. The
 parameter \(\beta_{2}\)
represents the
\(\textsc{e y}\) difference in the
\(\beta_{1}+\beta_{2}\) slopes.
\(\beta_{1}\)
 90000 95000 100000 10500
x
 Special Case: Nonlinear Functions of a Continuous Variable. In some applica-
tions, we expect the response to have some abrupt changes in behavior at certain
 values of an explanatory variable, even if the variable is continuous. For example,
 suppose that we are trying to model an individual's charitable contributions \(\left( y \right)\) in
terms of his or her wages (x). For 2007 data, a simple model we might entertain
 is given in Figure 3.9.
A rational for this model is that, in 2007, individuals paid 7.65\% of their
income for Social Security taxes up to \$97,500. No social security taxes are
excised on wages in excess of \$97,500. Thus, one theory is that, for wages in
excess of \$97,500, individuals have more disposal income per dollar and thus
 should be more willing to make charitable contributions.
To model this relationship, define the binary variable \(z\) to be zero if \(x <\)
\(9 7 , \! 5 0 0\) and to be one if \(x ~ \geq~ 9 7 , \! 5 0 0\) Define the regression function to be
E \(\cdot\ y \,=\, \beta_{0} \,+\, \beta_{1} x \,+\, \beta_{2} z ( x \,-\, 9 7 {,} 5 0 0 ) .\) This can be written as
\[
\textrm{E} y=\left\{\begin{matrix} {\beta_{0}+\beta_{1} x} & {x < 9 7 , 5 0 0} \\ {\beta_{0}-\beta_{2} ( 9 7 , \! 5 0 0 )+( \beta_{1}+\beta_{2} ) x} & {x \geq9 7 , \! 5 0 0} \\ \end{matrix} \right. .
\]

\clearpage
100 Multiple Linear Regression -- I
Figure 3.10 Plot of
expected commissions 240
(E y) versus number
of shares traded (x). 220
 The break at \(x \,=\, 1 0 0\) 200
reflects savings in
expenses. The lower E y 180 -
administrative
 slope for \(x \, \geq\, 1 0 0\) 160
reflects economies of
 scales in expenses. 140
120
60 80 100 120 140
x
To estimate this model, we would run a regression of \(y\) on two explanatory
variables, \(x_{1}=x\) and \(x_{2}=z \times( x \,-9 7 , 5 0 0 )\) . If \(\beta_{2} > 0\) then the marginal rate
 of charitable contributions is higher for incomes exceeding \$97,500.
Figure 3.9 illustrates this relationship, known as piecewise linear regression or
sometimes \(\mathrm{a}\) " brolen-stick" model. The sharp break in Figure 3.9 at \(x \,=9 7 , \! 5 0 0\)
 is called a kink. We have linear relationships above and below the kinks and have
used a binary variable to put the two pieces together. We are not restricted to one
 kink. For example, suppose that we want to do a historical study of federal taxable
income for 1992 single filers. Then, there were three tax brackets: the marginal
tax rate below \$21,450 was \(1 5 \not\! \gamma_{o}\) , above \$51,900 was \(3 1 {\cal V}_{o}\) , and in between was
28 percent. For this example, we would use two kinks, at 21,450 and 51,900.
Further, piecewise linear regression is not restricted to continuous response
functions. For example, suppose that we are studying the commissions paid to
stockbrokers (y) in terms of the number of shares purchased by a client (x).
We might expect to see the relationship illustrated in Figure 3.10. Here, the
discontinuity at \(x \,=\, 1 0 0\) reflects the administrative expenses of trading in odd
 lots, as trades of less than 100 shares are called. The lower marginal cost for trades
 in excess of 100 shares simply reflects the economies of scale for doing business
 in larger volumes. A regression model of this is E \(\boldsymbol{y}=\beta_{0} \,+\, \beta_{1} \boldsymbol{x} \,+\, \beta_{2} \boldsymbol{z} \,+\, \beta_{3} \boldsymbol{z} \boldsymbol{x}\) :
where \(z=0\) if \(x \, < \, 1 0 0\) and \(z=1\) if \(x \, \geq\, 1 0 0 .\) The regression function depicted
in Figure 3.10 is
\[
\textrm{E} y=\begin{cases} {\, \beta_{0} \,+\, \beta_{1} x_{1}} & {x \, < \, 1 0 0} \\ {\, \beta_{0} \,+\, \beta_{2} \,+\, ( \beta_{1} \,+\, \beta_{3} ) x_{1}} & {x \, \geq\, 1 0 0} \\ \end{cases} .
\]
3.6 Further Reading and References
For proofs of the Chapter 3 results, see Goldberger (1991). Nonlinear regression
models are discussed in, for example, Bates and Watts (1988).
 Chapter 3 has introduced the fundamentals of multiple linear regression. Chap-
ter \(4\) will widen the scope by introducing categorical variables and statistical
\clearpage
3.7 Exercises 101
inference methods for handling several coefficients simultaneously. Chapter \(5\)
will introduce techniques to help you pick appropriate variables in a multiple
 linear regression model. Chapter \(6\) is a synthesis chapter, discussing model inter-
 pretation, variable selection, and data collection.
Chapter References
 Bates, Douglas M., and D. G. Watts (1988). Nonlinear Regression Analysis and its Applications.
John Wiley and Sons, New York.
Lemaire, Jean (2002). Why do females live longer than males? North American Actuarial
 Journal 6, no. 4, 21- 37.
Goldberger, Arthur (1991). A Course in Econometrics. Harvard University Press, Cambridge,
Massachusetts.
Plackett, \(\mathrm{R}\) . L. (1960). Regression Analysis. Clarendon Press, Oxford, U.K.
 Segal, Dan (2002). An economic analysis of life insurance company expenses. North American
 Actuarial Journal \(6\) , no. \(4\) 81- 94.
3.7 Exercises
3.1. Consider a fictitious dataset of \(n \,=\, 1 0 0\) observations with \(s_{y}=8 0 .\) We run
 a regression with three explanatory variables to get \(s \,=5 0\) .
a. Calculate the adjusted coefficient of determination, \(R_{a}^{2}\) .
b. Complete the ANOVA table.
ANOVA Table
Source Sum of Squares \(d f\) Mean Square
Regression
Error
Total
\(\mathrm{c}\) .Calculate the (unadjusted) coefficient of determination, \(R^{2}\)
3.2. Consider a fictitious dataset of \(n \,=\, 1 0 0\) observations with \(s_{y}=8 0 .\) We run
 a regression with three explanatory variables to get \(s \,=5 0\) We also get
\[
\left( \mathbf{X}^{\prime} \mathbf{X} \right)^{-1}=\left( \begin{matrix} {1 0 0} & {2 0} & {2 0} & {2 0} \\ {2 0} & {9 0} & {3 0} & {4 0} \\ {2 0} & {3 0} & {8 0} & {5 0} \\ {2 0} & {4 0} & {5 0} & {7 0} \\ \end{matrix} \right) .
\]
 a. Determine the standard error of \(b_{3} , \, s e ( b_{3} )\) .
b. Determine the estimated covariance between \(b_{2}\) and \(b_{3}\) .
 c. Determine the estimated correlation between \(b_{2}\) and \(b_{3}\) .
d. Determine the estimated variance of \(4 b_{2}+3 b_{3}\) .
\clearpage
102 Multiple Linear Regression -- I
3.3. Consider the following small fictitious dataset. You will be fitting a regres-
 sion model to \(y\) using two explanatory variables, \(x_{1}\) and \(x_{2}\) .
\[
\begin{array} {l l l l l} \hline{} & {} & {} & {} & {} \\ {i} & {1} & {2} & {3} & {4} \\ {x_{i , 1}} & {-1} & {2} & {4} & {6} \\ {x_{i , 2}} & {0} & {0} & {1} & {1} \\ {y_{i}} & {0} & {1} & {5} & {8} \\ \hline{} \\ \end{array}
\]
From the fitted regression model, we have \(s \,=1 . 3 7 3\) and
\[
\mathbf{b}=\left( \begin{matrix} {0 . 1 5} \\ {0 . 6 9 2} \\ {2 . 8 8} \\ \end{matrix} \right) \mathrm{~ a n d ~} \ \left( \mathbf{X}^{*} \mathbf{X} \right)^{-1}=\left( \begin{matrix} {0 . 5 3 8 4 6} & {-0 . 0 7 6 9 2} & {-0 . 1 5 3 8 5} \\ {-0 . 0 7 6 9 2} & {0 . 1 5 3 8 5} & {-0 . 6 9 2 3 1} \\ {-0 . 1 5 3 8 5} & {-0 . 6 9 2 3 1} & {4 . 1 1 5 3 8} \\ \end{matrix} \right) .
\]
a. Write down the vector of dependent variables, \({\bf y}\) , and the matrix of
 explanatory variables, \(\mathbf{X}\) .
Determine the numerical value for \(\widehat{y}_{3}\) , the fitted value for the third
b.
observation.
 c. Determine the numerical value for \(s e ( b_{2} )\) .
@ EMPIRICAL d. Determine the numerical value for \(t ( b_{1} )\) .
Filename is 3.4. Wisconsin Lottery. Section 2.1 described a sample of \(n=5 0\) geographic
"WiscLottery" areas (Zip codes) containing sales data on the Wisconsin state lottery ( \(y=\)
 SALES). In that section, sales were analyzed using a basic linear regression
 model with \(x \,=\, P \, O \, P\) , the area population, as the explanatory variable.
 This exercise extends that analysis by introducing additional explanatory
 variables given in Table 3.11.
Table 3.11 Lottery,
Economic, and Lottery Characteristics
Demographic
Characteristics of 50 SALES Online Lottery Sales to Individual Consumers
Wisconsin Zip Codes
Economic and demographic characteristics
PERPERHH Persons per household
MEDSCHYR Median years of schooling
MEDHVL Median home value in \$1000s for owner-occupied homes
PRCRENT Percentage of housing that is renter occupied
PRC55P Percentage of population that is 55 or older
HHMEDAGE Household median age
MEDINC Estimated median household income, in \$1000s
POP Population, in thousands
a. Produce a table of summary statistics for all variables. One Zip code
(observation 11, \(\mathrm{Z i p} \,=\, 5 3 2 1 1\) , Shorewood, Wisconsin, a suburb of
Milwaukee) appears to have unusually large values of MEDSCHYR
 and MEDHVL. For this observation, how many standard deviations is
the value of MEDSCHYR above the mean? For this observation, how
 many standard deviations is the value of MEDHVL above the mean?
\clearpage
3.7 Exercises 103
b. Produce a table of correlations. What three variables are most highly
correlated with SALES?
 c. Produce a scatterplot matrix of all explanatory variables and SALES.
In the plot of MEDSCHYR versus SALES, describe the position of
observation 11.
 d. Fit a linear model of SALES on all eight explanatory variables. Sum-
 marize the fit of this model by citing the residual standard deviation, \(s\) :
 the coefficient of determination, \(R^{2}\) ; and its adjusted version, \(R_{a}^{2}\) .
e. Based on your part (d) model fit, is MEDSCHYR a statistically sig
nificant variable? To respond to this question, use a formal test of
 hypothesis. State your null and alternative hypotheses, decision-making
criterion, and decision-making rule.
f. Now fit a more parsimonious model, using SALES as the dependent
variable and MEDSCHYR, MEDHVL, and POP as explanatory vari-
ables. Summarize the fit of this model by citing the residual standard
deviation, \(s\) ; the coefficient of determination, \(R^{2}\) ; and its adjusted ver-
sion, \(R_{a}^{2}\) . How do these values compare to the model fit in part (d)?
g. Note that the sign of the regression coefficient associated with MED-
SCHYR is now negative. To help interpret this coefficient, compute the
corresponding partial correlation coefficient. What is the interpretation
of this coefficient?
h. To get further insights into the relation between MEDSCHYR and
SALES, produce an added variable plot controlling for the effects of
MEDHVL and POP. Check that the correlation associated with this
 plot agrees with your answer in part \(\mathrm{( g )}\) .
i. Rerun the regression in part (f), after removing observation 11. Cite
the basic summary statistics from this regression. For this model fit,
is MEDSCHYR a statistically significant variable? To respond to this
 question, use a formal test of hypothesis. State your null and alternative
hypotheses, decision-making criterion, and decision-making rule.
 Rerun the regression in part (f), after removing observation \(9\) . Cite the
j.
 basic summary statistics from this regression. @ EMPIRICAL
3.5. Insurance Company Expenses. This exercise considers insurance com- Filename is
 pany data from the NAIC and described in Exercise 1.6. "NAICExpense"
Table 3.12 describes several variables that can be used to explain
expenses. As with Segal's (2002) study of life insurers, firm " outputs"
consist of premiums written (for property and casualty, these are sub-
divided into personal and commercial lines) and losses (subdivided into
short and long tail lines). ASSETS and CASH are commonly used mea-
sures of the size of a company. GROUP, STOCK, and MUTUAL describe
 the organizational structure. Firm " inputs"were gathered from the Bureau
 of Labor Statistics (BLS, from the Occupational Employee Statistics pro-
 gram). WAGESTAFF is calculated as the average wage in the state where
the insurance company is headquartered. AGENTWAGE is calculated as
\clearpage
104 Multiple Linear Regression -- I
 the weighted average of annual wages of the brokerage industry, weighted
 by the percentage of gross premium written in each state.
Table 3.12 Insurer
 Expense Variables Variable Description
 NAIC Variables
EXPENSES Total expenses incurred, in millions of dollars
LOSSLONG Losses incurred for long tail lines, in millions of dollars
LOSSSHORT Losses incurred for short tail lines, in millions of dollars
GPWPERSONAL Gross premium written for personal lines, in millions of dollars
GPWCOMM Gross premium written for commercial lines, in millions of dollars
ASSETS Net admitted assets, in millions of dollars
CASH Cash and invested assets, in millions of dollars
GROUP Indicates whether the company is affiliated
 STOCK Indicates whether the company is a stock company
 MUTUAL indicates whether the company is a mutual company
 BLS Variables
STAFFWAGE Annual average wage of the insurer's administrative staff,
 in thousands of dollars
AGENTWAGE Annual average wage of the insurance agent, in thousands of dollars
 A preliminary inspection of the data showed that many firms did not
report any insurance losses incurred in 2005. For this exercise, we consider
the 384 companies with some losses in the file " NAICExpense."
 a. Produce summary statistics of the response variable and the (nonbinary)
 explanatory variables. Note the pattern of skewness for each variable.
 Note that many variables have negative values.
b. Transform each nonbinary variable through the modified logarithm
transform, \(\operatorname{l n} ( 1+x )\) Produce summary statistics of these modi-
fied nonbinary explanatory variables. Let LNEXPENSES \((=\operatorname{l n} ( 1+\)
EXPENSES)) denote the modified expense variable.
For subsequent analysis, use only the modified variables described
in part (b).
c. Produce a table of correlations for the nonbinary variables. What three
Variables are most highly correlated with LNEXPENSES?
d. Provide a boxplot of LNEXPENSES by level of GROUP. Which level
of group has higher expenses?
e. Fit a linear model of LNEXPENSES on all eleven explanatory vari
 ables. Summarize the fit of this model by citing the residual standard
 deviation, \(s\) ; the coefficient of determination, \(R^{2}\) ; and its adjusted ver-
 sion, \(R_{a}^{2}\) .
f. Fit a linear model of LNEXPENSES on a reduced model using eight
explanatory variables, dropping CASH, STOCK, and MUTUAL. For
the explanatory variables, include assets, GROUP, both versions of
losses and gross premiums, as well as the two BLSvariables.
\clearpage
3.7 Exercises 105
f(i). Summarize the fit of this model by citing \(s , \, \, R^{2}\) and \(R_{a}^{2}\)
f(ii). Interpret the coefficient associated with commercial lines gross
 premiums on the logarithmic scale.
f(iii). Suppose that GPWCOMM increases by \$1.00, how much do we
expect EXPENSES to increase? Use your answer in part f(ii)
and median values of GPWCOMM and EXPENSES for this
question.
g. Square each of the two loss and the two gross premium variables. Fit
 a linear model of LNEXPENSES on a reduced model using twelve
explanatory variables, the eight variables in part (f), and the four addi-
 tional squared terms just created.
\(\mathrm{g ( i )}\) .Summarize the fit of this model by citing \(s , \; R^{2}\) , and \(R_{a}^{2}\) .
\(\mathrm{g ( i i )}\) . Do the quadratic variables appear to be useful explanatory
variables?
h. Now omit the two BLS variables, so you are fitting a model of LNEX-
PENSES on assets, GROUP, both versions of losses and gross premi-
ums, as well as quadratic terms. Summarize the fit of this model by
citing \(s , \; R^{2}\) , and \(R_{a}^{2}\) Comment on the number of observations used to
fit this model compared to part (f).
i. Drop the quadratic terms in part (g) and add interaction terms with
 the dummy variable GROUP. Thus, there are now 11 variables, assets,
GROUP, both versions of losses and gross premiums, as well as inter-
actions of GROUP with assets and both versions of losses and gross
premiums.
i(i). Summarize the fit of this model by citing \(s , \; R^{2}\) , and \(R_{a}^{2}\) .
i(ii). Suppose that GPWCOMM increases by \(\S1 . 0 0\) , how much do
we expect EXPENSES to increase for \(\mathrm{G R O U P} ~=~ 0\) compa-
nies? Use the median values of GPWCOMM and EXPENSES of
\(\mathrm{G R O U P}=0\) companies for this question.
i(ii). Suppose that GPWCOMM increases by \$1.00, how much do
we expect EXPENSES to increase for GROUP \(= 1\) compa-
nies? Use the median values of GPWCOMM and EXPENSES of
\(\mathrm{G R O U P}=1\) companies for this question. @ EMPIRICAL
3.6. National Life Expectancies. We continue the analysis begun in ExercisesFilename is
1.7 and 2.22. Now fit a regression model on LIFEEXP using three explana-"UNLifeExpectancy"
tory variables, FERTILITY, PUBLICEDUCATION, and lnHEALTH (the
 natural logarithmic transform of PRIVATEHEALTH).
 a. Interpret the regression coefficient associated with public education.
b. Interpret the regression coefficient associated with health expenditures
without using the logarithmic scale for expenditures.
c.Based on the model fit, is PUBLICEDUCATION a statistically sig-
nificant variable? To respond to this question, use a formal test of
 hypothesis. State your null and alternative hypotheses, decision-making
criterion, and decision-making rule.
\clearpage
106 Multiple Linear Regression -- I
Figure 3.11 Added 2 : S : \(\aleph_{\circ}\) 0 0
variable plot of PUB- 5 S
lnHEALTH. TE d g0
LICEDUCATION
versus LIFEEXP, 5 0 0
6
 controlling for 00
FERTILITY and
2 0
2
4 2 0 2 4 6 8
residuals(PUBLICEDUCATION)
 d. The negative sign of the PUBLICEDUCATION coefficient is surpris-
 ing, in that the sign of the correlation between PUBLICEDUCATION
 and LIFEEXP is positive and intuition suggests a positive relation. To
check this result, an added variable plot appears in Figure 3.11.
\(\mathrm{d ( i )}\) .For an added variable plot, describe its purpose and a method for
 producing it.
 d(ii). Calculate the correlation corresponding to the added variable plot
 that appears in Figure 3.11.
\clearpage

\end{document}
